% Lesson 5 — Grammar: Review tenses in the previous years + present perfect continuous and future continuous tenses — Anglais Tle A — Cours signé OnBuch+
% Programme officiel MINESEC. Compiler avec : tectonic EN05-grammar-review-tenses-in-the-previous-years-present-per.tex
\newcommand{\DOCMATIERE}{Anglais}
\newcommand{\DOCNIVEAU}{Tle A}
\newcommand{\DOCMODULE}{Chapter 1 : Family and Social Life}
\newcommand{\DOCLECON}{EN05}
\newcommand{\DOCTITRE}{Lesson 5 — Grammar: Review tenses in the previous years + present perfect continuous and future continuous tenses}
% OnBuch+ — préambule « cours signés » ENRICHI (compilé avec tectonic / XeLaTeX)
% Système visuel « Pop » d'OnBuch+ : fond crème (jamais blanc), bordures encre
% épaisses, ombres « dures » décalées, titres Archivo Black en capitales.
% Version MATHÉMATIQUES : graphes (pgfplots/tikz), boîtes pédagogiques
% (définition, propriété, méthode, attention, le savais-tu, exercice résolu,
% exercice, TP, bilan), page de garde et signature OnBuch+ ancrée en bas.
%
% Macros attendues dans main.tex AVANT \input{preamble} :
%   \DOCMATIERE  \DOCNIVEAU  \DOCMODULE  \DOCLECON (numéro)  \DOCTITRE
\documentclass[11pt]{article}

\usepackage{fontspec}
\usepackage[english,french]{babel} % français = langue principale ; l'anglais via \eng{..} / textebox
\usepackage[a4paper,margin=2cm,top=3.4cm,bottom=2.8cm,headheight=1.4cm,footskip=1.3cm]{geometry}
\usepackage{amsmath,amssymb,amsfonts}
\usepackage{mathtools} % \xmapsto, \coloneqq... souvent utilisés par les modèles
\usepackage{cancel} % simplifications barrées (bilans de forces)
% \abs{z} (module, valeur absolue) et \norm{u} (norme), utilisés par les modèles
\providecommand{\abs}{}\let\abs\relax\DeclarePairedDelimiter{\abs}{\lvert}{\rvert}
\providecommand{\norm}{}\let\norm\relax\DeclarePairedDelimiter{\norm}{\lVert}{\rVert}
\usepackage[table]{xcolor} % [table] : fournit aussi \rowcolors (alternance de lignes)
\usepackage{pagecolor}
\usepackage{tikz}
\usetikzlibrary{arrows.meta,positioning,calc,shapes.geometric,shapes.misc,shapes.arrows,decorations.pathmorphing,decorations.pathreplacing,decorations.markings,patterns,shadows,mindmap,trees,fit,backgrounds,matrix,intersections,angles,quotes}
\usepackage{pgfplots}
\usepackage[version=4]{mhchem} % équations nucléaires \ce{^{235}_{92}U}
\usepackage[european,siunitx]{circuitikz} % schémas électriques
\pgfplotsset{compat=1.18}
\usepgfplotslibrary{fillbetween,groupplots} % aires entre courbes (calcul intégral)
\usepackage{enumitem}
\usepackage{fancyhdr}
% [most] charge toutes les bibliothèques tcolorbox (listings, minted, raster,
% documentation...) et gonfle inutilement la mémoire TeX sur un document long
% (~300 boîtes) : on ne charge que ce qui est réellement utilisé.
\usepackage[skins,breakable]{tcolorbox}
\usepackage{graphicx}
\usepackage{colortbl}
\usepackage{booktabs}
\usepackage{tabularx}
\usepackage{diagbox} % case d'en-tête diagonale des tableaux à double entrée
% Colonnes extensibles centrée / gauche / droite (C, L, R), souvent utilisées par les modèles
\newcolumntype{C}{>{\centering\arraybackslash}X}
\newcolumntype{L}{>{\raggedright\arraybackslash}X}
\newcolumntype{R}{>{\raggedleft\arraybackslash}X}
\usepackage{array}
\usepackage{multicol}
\usepackage{multirow}
\usepackage{siunitx}
% parse-numbers=false : accepte aussi \SI{2,5 \times 10^{-3}}{...}
\sisetup{locale=FR,output-decimal-marker={,},group-minimum-digits=4,parse-numbers=false}
\DeclareSIUnit\ohm{\ensuremath{\symup{\Omega}}} % U+2126 absent de la police
\DeclareSIUnit\division{div} % réglages d'oscilloscope (ms/div, V/div)
\DeclareSIUnit\tr{tr} % tours (vitesse de rotation, ex. \SI{1500}{\tr\per\minute})
\DeclareSIUnit\molar{M} % concentration molaire (ex. \SI{3}{\molar}, \SI{1}{\micro\molar})
\DeclareSIUnit\an{an} % année (le modèle confond parfois avec \year, un compteur TeX) — ex. \SI{5}{\kilo\meter\per\an}
\DeclareSIUnit\FCFA{FCFA} % franc CFA (ex. \SI{500}{\FCFA\per\kilo\gram})
\DeclareSIUnit\degreeBrix{\textdegree Brix} % teneur en sucre d'un sirop/jus (réfractométrie)
\DeclareSIUnit\month{mois} % le modèle utilise parfois \month, un compteur TeX, comme unité de durée
\usepackage{qrcode}
% \degree : commande fréquemment utilisée par les modèles pour un angle ou une
% température en dehors de \SI{}{} (non fournie par les paquets chargés).
\providecommand{\degree}{^{\circ}}
% \qed (fin de démonstration), utilisé par les modèles sans amsthm
\providecommand{\qed}{\hfill$\square$}
% \barre : double barre des valeurs interdites dans un tableau de variations (tkz-tab)
\providecommand{\barre}{\ensuremath{\|}}
% environnement proof (amsthm non chargé) utilisé par les modèles
\ifdefined\proof\else\newenvironment{proof}[1][Démonstration]{\par\noindent\textit{#1.}\ }{\qed\par}\fi
\usepackage{lastpage}
\usepackage{hyperref}
\hypersetup{hidelinks}

% ============================================================
% Palette « Pop » OnBuch+ — mêmes hex que la classe P (plus_ui.dart)
% ============================================================
\definecolor{poporange}{HTML}{F59321}
\definecolor{popdark}{HTML}{E07A0C}
\definecolor{popcream}{HTML}{FFF6E8}
\definecolor{popink}{HTML}{1C1714}
\definecolor{poppurple}{HTML}{7A5AE0}
\definecolor{popgreen}{HTML}{2E9E6A}
\definecolor{poppink}{HTML}{E0457B}
\definecolor{popblue}{HTML}{2F7DE1}
\definecolor{popgold}{HTML}{C98A12}
\definecolor{popmuted}{HTML}{8A7F76}
\definecolor{popline}{HTML}{EADFD2}
\definecolor{popgray}{HTML}{8A7F76}
\colorlet{popgrey}{popgray}
% Alias tolérants (le modèle nomme parfois les couleurs différemment)
\colorlet{popwhite}{white}
\colorlet{popblack}{popink}
\colorlet{popred}{poppink}
\colorlet{popyellow}{popgold}
% Teintes claires (fonds de tableaux/encadrés) — le modèle nomme parfois
% les couleurs avec un suffixe L (ex. popblueL) sans les avoir définies.
\colorlet{poporangeL}{poporange!20}
\colorlet{popdarkL}{popdark!20}
\colorlet{poppurpleL}{poppurple!20}
\colorlet{popgreenL}{popgreen!20}
\colorlet{poppinkL}{poppink!20}
\colorlet{popblueL}{popblue!20}
\colorlet{popgoldL}{popgold!20}
\colorlet{popredL}{popred!20}
\colorlet{popgrayL}{popgray!20}
% Teintes claires pour fonds de tableaux / zones
\colorlet{popblueL}{popblue!10!white}
\colorlet{poporangeL}{poporange!14!white}
\colorlet{popgreenL}{popgreen!12!white}
\colorlet{poppurpleL}{poppurple!12!white}
\colorlet{poppinkL}{poppink!12!white}
\colorlet{popgoldL}{popgold!14!white}

\pagecolor{popcream}

% ============================================================
% Typographie
% ============================================================
\defaultfontfeatures{Ligatures=TeX}
\setmainfont{PlusJakartaSans-Medium.ttf}[Path=../../../fonts/,BoldFont=PlusJakartaSans-Bold.ttf]
% Maths en sans-serif (Fira Math) avec chiffres et lettres droites de Plus
% Jakarta, pour que les formules se fondent dans le texte.
\usepackage[math-style=ISO,bold-style=ISO]{unicode-math}
\setmathfont{FiraMath-Regular.otf}[Path=../../../fonts/]
\setmathfont{PlusJakartaSans-Medium.ttf}[Path=../../../fonts/,range={up/num,up/{Latin,latin},\mathup}]
\usepackage{icomma} % virgule décimale sans espace en mode math
% Caractères mathématiques Unicode absents des polices de texte (Plus Jakarta,
% Archivo Black) : rendus par leur équivalent mathématique, en texte comme en
% formule — sinon ils s'affichent en carré vide (ex. « Arithmétique dans ℤ »).
\usepackage{newunicodechar}
\newunicodechar{ℝ}{\ensuremath{\mathbb{R}}}
\newunicodechar{ℤ}{\ensuremath{\mathbb{Z}}}
\newunicodechar{ℂ}{\ensuremath{\mathbb{C}}}
\newunicodechar{ℕ}{\ensuremath{\mathbb{N}}}
\newunicodechar{ℚ}{\ensuremath{\mathbb{Q}}}
\newunicodechar{𝔻}{\ensuremath{\mathbb{D}}}
\newunicodechar{α}{\ensuremath{\alpha}}
\newunicodechar{β}{\ensuremath{\beta}}
\newunicodechar{γ}{\ensuremath{\gamma}}
\newunicodechar{δ}{\ensuremath{\delta}}
\newunicodechar{ε}{\ensuremath{\varepsilon}}
\newunicodechar{θ}{\ensuremath{\theta}}
\newunicodechar{λ}{\ensuremath{\lambda}}
\newunicodechar{μ}{\ensuremath{\mu}}
\newunicodechar{σ}{\ensuremath{\sigma}}
\newunicodechar{τ}{\ensuremath{\tau}}
\newunicodechar{φ}{\ensuremath{\varphi}}
\newunicodechar{ψ}{\ensuremath{\psi}}
\newunicodechar{ω}{\ensuremath{\omega}}
\newunicodechar{Σ}{\ensuremath{\Sigma}}
% majuscules produites par \MakeUppercase dans les titres (α-aminés -> Α)
\newunicodechar{Α}{\ensuremath{\alpha}}
\newunicodechar{Β}{\ensuremath{\beta}}
\newunicodechar{Γ}{\ensuremath{\Gamma}}
\newunicodechar{Δ}{\ensuremath{\Delta}}
\newunicodechar{Ε}{\ensuremath{\varepsilon}}
\newunicodechar{Θ}{\ensuremath{\Theta}}
\newunicodechar{Λ}{\ensuremath{\Lambda}}
\newunicodechar{Μ}{\ensuremath{\mu}}
\newunicodechar{Τ}{\ensuremath{\tau}}
\newunicodechar{Φ}{\ensuremath{\Phi}}
\newunicodechar{Ψ}{\ensuremath{\Psi}}
\newunicodechar{Ω}{\ensuremath{\Omega}}
\newunicodechar{↦}{\ensuremath{\mapsto}}
\newunicodechar{∈}{\ensuremath{\in}}
\newunicodechar{∉}{\ensuremath{\notin}}
\newunicodechar{∧}{\ensuremath{\wedge}}
\newunicodechar{⇔}{\ensuremath{\Leftrightarrow}}
\newunicodechar{⟺}{\ensuremath{\Longleftrightarrow}}
\newunicodechar{⇒}{\ensuremath{\Rightarrow}}
\newunicodechar{⟹}{\ensuremath{\Longrightarrow}}
\newunicodechar{∘}{\ensuremath{\circ}}
\newunicodechar{∪}{\ensuremath{\cup}}
\newunicodechar{∩}{\ensuremath{\cap}}
\newunicodechar{≡}{\ensuremath{\equiv}}
\newunicodechar{⊂}{\ensuremath{\subset}}
\newunicodechar{⊥}{\ensuremath{\perp}}
\newunicodechar{∃}{\ensuremath{\exists}}
\newunicodechar{∀}{\ensuremath{\forall}}
\newunicodechar{∛}{\ensuremath{\sqrt[3]{\,}}}
\newunicodechar{∜}{\ensuremath{\sqrt[4]{\,}}}
\newunicodechar{⃗}{\ensuremath{{}^{\to}}}
\newunicodechar{ⁿ}{\ifmmode^{n}\else\textsuperscript{\ensuremath{n}}\fi}
\newunicodechar{ˣ}{\ifmmode^{x}\else\textsuperscript{\ensuremath{x}}\fi}
\newunicodechar{ᵇ}{\ifmmode^{b}\else\textsuperscript{\ensuremath{b}}\fi}
\newunicodechar{ʳ}{\ifmmode^{r}\else\textsuperscript{\ensuremath{r}}\fi}
\newunicodechar{ᵘ}{\ifmmode^{u}\else\textsuperscript{\ensuremath{u}}\fi}
\newunicodechar{ʸ}{\ifmmode^{y}\else\textsuperscript{\ensuremath{y}}\fi}
\newunicodechar{ᵃ}{\ifmmode^{a}\else\textsuperscript{\ensuremath{a}}\fi}
\newunicodechar{ᵉ}{\ifmmode^{e}\else\textsuperscript{\ensuremath{e}}\fi}
\newunicodechar{ᵏ}{\ifmmode^{k}\else\textsuperscript{\ensuremath{k}}\fi}
\newunicodechar{ᵖ}{\ifmmode^{p}\else\textsuperscript{\ensuremath{p}}\fi}
\newunicodechar{ᴺ}{\ifmmode^{N}\else\textsuperscript{\ensuremath{N}}\fi}
\newunicodechar{ᶜ}{\ifmmode^{c}\else\textsuperscript{\ensuremath{c}}\fi}
\newunicodechar{ʰ}{\ifmmode^{h}\else\textsuperscript{\ensuremath{h}}\fi}
\newunicodechar{ᵠ}{\ifmmode^{q}\else\textsuperscript{\ensuremath{q}}\fi}
\newunicodechar{ₙ}{\ifmmode_{n}\else\textsubscript{\ensuremath{n}}\fi}
\newunicodechar{ₐ}{\ifmmode_{a}\else\textsubscript{\ensuremath{a}}\fi}
\newunicodechar{ᵢ}{\ifmmode_{i}\else\textsubscript{\ensuremath{i}}\fi}
\newunicodechar{ᵣ}{\ifmmode_{r}\else\textsubscript{\ensuremath{r}}\fi}
\newunicodechar{ₖ}{\ifmmode_{k}\else\textsubscript{\ensuremath{k}}\fi}
\newunicodechar{ᵦ}{\ifmmode_{\beta}\else\textsubscript{\ensuremath{\beta}}\fi}
\newunicodechar{ₓ}{\ifmmode_{x}\else\textsubscript{\ensuremath{x}}\fi}
\newfontfamily\popdisplay{ArchivoBlack-Regular.ttf}[Path=../../../fonts/]
\newfontfamily\poplogo{SpaceGrotesk-Bold.ttf}[Path=../../../fonts/]
\newfontfamily\popsemi{PlusJakartaSans-SemiBold.ttf}[Path=../../../fonts/]
\newfontfamily\popxbold{PlusJakartaSans-ExtraBold.ttf}[Path=../../../fonts/]
\newfontfamily\popxboldls{PlusJakartaSans-ExtraBold.ttf}[Path=../../../fonts/,LetterSpace=3.0]

\setlist[enumerate]{leftmargin=1.7em,itemsep=5pt,topsep=4pt}
\setlist[itemize]{leftmargin=1.7em,itemsep=4pt,topsep=4pt,label={\color{poporange}\textbullet}}
\setlength{\parindent}{0pt}
\setlength{\parskip}{5pt}
\renewcommand{\arraystretch}{1.25}

% pgfplots : style Pop par défaut
\pgfplotsset{
  every axis/.append style={axis line style={popink,line width=1pt},tick style={popink},
    grid style={popline,line width=0.5pt},label style={font=\small},tick label style={font=\footnotesize}},
  popcurve/.style={poporange,line width=1.8pt,no markers,smooth},
}
% Même style disponible dans un \draw TikZ simple (hors pgfplots)
\tikzset{popcurve/.style={poporange,line width=1.8pt}}
\tikzset{popgrid/.style={popline,very thin}}
\colorlet{popcurve}{poporange} % le nom du style est parfois utilisé comme couleur

% ============================================================
% Pill / squiggle
% ============================================================
\newcommand{\poppill}[3]{%
  \tikz[baseline=-0.55ex]\node[draw=popink,line width=1.2pt,rounded corners=2.6pt,%
    fill=#2,inner xsep=6.5pt,inner ysep=3pt]{%
    \popxboldls\fontsize{7}{7}\selectfont\color{#3}\MakeUppercase{#1}};%
}
\newcommand{\popsquiggle}[1]{%
  \tikz[baseline=-0.4ex]{
    \draw[#1,line width=2.6pt,line cap=round]
      plot [smooth] coordinates {(0,0.09) (0.34,0.01) (0.68,0.21) (1.02,0.01) (1.36,0.10)};
  }%
}
% Mot-clé surligné (marqueur orange sous le texte)
\newcommand{\cle}[1]{{\popxbold\color{popdark}#1}}

% ============================================================
% En-tête / pied
% ============================================================
\pagestyle{fancy}
\fancyhf{}
\renewcommand{\headrulewidth}{0pt}
\renewcommand{\footrulewidth}{0pt}
\lhead{\raisebox{-0.15ex}{\poplogo\fontsize{13}{13}\selectfont\color{popink}OnBuch\color{poporange}+}}
\rhead{\poppill{\DOCMATIERE\ \textbullet\ \DOCNIVEAU\ \textbullet\ Leçon \DOCLECON}{white}{popink}}
\lfoot{\popxbold\fontsize{7.6}{7.6}\selectfont\color{popmuted}\MakeUppercase{Cours signé OnBuch+}\ \textbullet\ {\popsemi\DOCTITRE}}
\rfoot{\tikz[baseline=-0.6ex]\node[rounded corners=8.5pt,fill=popink,minimum height=17pt,minimum width=17pt,inner xsep=5pt,inner sep=0]{\popxbold\fontsize{8}{8}\selectfont\color{popcream}\thepage\,/\,\pageref*{LastPage}};}

% ============================================================
% Titres
% ============================================================
\newcommand{\chaptitre}[2]{%
  \begin{tcolorbox}[enhanced,colback=poporange,colframe=popink,boxrule=2.6pt,arc=11pt,
      drop shadow=popink,left=15pt,right=15pt,top=12pt,bottom=11pt,before skip=3pt,after skip=18pt]
    \poppill{#2}{popink}{popcream}\par\vspace{9pt}
    {\raggedright\hyphenpenalty=10000\exhyphenpenalty=10000\popdisplay\fontsize{22}{24}\selectfont\color{popcream}\MakeUppercase{#1}\par}
  \end{tcolorbox}
}
% Section numérotée : pastille orange avec le numéro + titre Archivo
\newcounter{popsec}
\newcommand{\coursec}[1]{%
  \stepcounter{popsec}\setcounter{popsub}{0}%
  \par\vspace{16pt}\noindent
  \begin{minipage}{\linewidth}
  \tikz[baseline=-0.7ex]\node[draw=popink,line width=1.6pt,fill=poporange,rounded corners=4pt,minimum size=22pt,inner sep=0]{\popdisplay\fontsize{12}{12}\selectfont\color{popcream}\thepopsec};%
  \hspace{8pt}{\popdisplay\fontsize{15}{17}\selectfont\color{popink}\MakeUppercase{#1}}\par
  \vspace{4pt}\noindent{\color{poporange}\rule{36pt}{3.6pt}}
  \end{minipage}\par\nopagebreak\vspace{8pt}
}
\newcounter{popsub}
\newcommand{\courssub}[1]{%
  \stepcounter{popsub}%
  \par\vspace{9pt}\noindent{\popxbold\fontsize{11.5}{13}\selectfont\color{popink}{\color{poporange}\thepopsec.\thepopsub}\hspace{6pt}#1}\par\nopagebreak\vspace{4pt}
}

% ============================================================
% Boîtes pédagogiques — même recette : fond blanc, bordure épaisse colorée,
% ombre dure encre, badge plein. Titre optionnel [#1].
% ============================================================
\tcbset{popbase/.style={breakable,enhanced,colback=white,boxrule=2.3pt,arc=9pt,
  shadow={3pt}{-3pt}{0pt}{popink},left=14pt,right=14pt,top=11pt,bottom=11pt,before skip=11pt,after skip=15pt,
  fontupper=\popsemi\fontsize{10.6}{14.5}\selectfont\color{popink},
  fontlower=\popsemi\fontsize{10.6}{14.5}\selectfont\color{popink}}}
% Coche dessinée (le glyphe ✓ n'existe pas dans les polices du document)
\newcommand{\popcheck}[1]{\tikz[baseline=-0.5ex]\draw[#1,line width=1.6pt,line cap=round,line join=round](-0.13,0.0)--(-0.04,-0.09)--(0.13,0.1);}
\providecommand{\coche}{\popcheck{popgreen}}
\providecommand{\resultat}[1]{\par\smallskip\noindent\fcolorbox{popgreen}{popgreenL}{\parbox{\dimexpr\linewidth-2\fboxsep-2\fboxrule\relax}{\textbf{Résultat.} #1}}\par\smallskip}
\newcommand{\popboxhead}[3]{\poppill{#1}{#2}{white}\ifx\relax#3\relax\else\hspace{6pt}{\popxbold\fontsize{10.6}{12}\selectfont\color{popink}#3}\fi\par\vspace{7pt}}

\newtcolorbox{aretenir}[1][]{popbase,colframe=popgreen,before upper={\popboxhead{À retenir}{popgreen}{#1}}}
% Texte en anglais (document de lecture, dialogue, extrait) : cadre
% d'étude. Un titre optionnel (source, auteur) s'affiche dans l'en-tête.
\newtcolorbox{textebox}[1][]{popbase,colframe=popblue,before upper={\popboxhead{Text}{popblue}{#1}\begin{otherlanguage}{english}},after upper={\end{otherlanguage}}}
% Marques ✓ / ✗ (forme correcte / fautive) souvent utilisées par les modèles
\providecommand{\cross}{\textcolor{poppink}{\ensuremath{\times}}}
\providecommand{\xmark}{\cross}\providecommand{\crossmark}{\cross}\providecommand{\cmark}{\ensuremath{\checkmark}}
% Ligne de réponse / justification à compléter (inventée par les modèles)
\providecommand{\justification}[1]{\par\noindent\textit{Justification~:} \dotfill\par}
% \textipa (package tipa non chargé) : la police ne couvre pas l'API -> simple passage
\providecommand{\textipa}[1]{#1}
% Soulignement (ulem non chargé) : \uline, \ul
\providecommand{\uline}[1]{\underline{#1}}\providecommand{\ul}[1]{\underline{#1}}
% \glossaire{\item ...} inventée par les modèles : liste de glossaire
\providecommand{\glossaire}[1]{\begin{itemize}#1\end{itemize}}
% \alert (beamer) inventée par les modèles : texte mis en évidence
\providecommand{\alert}[1]{\textbf{\textcolor{poppink}{#1}}}
% variantes ASCII d'ordinaux écrites en macro
\providecommand{\iere}{\textsuperscript{re}}\providecommand{\ieme}{\textsuperscript{e}}\providecommand{\ere}{\textsuperscript{re}}\providecommand{\eme}{\textsuperscript{e}}\providecommand{\er}{\textsuperscript{er}}\providecommand{\ie}{\textsuperscript{e}}
% Ordinaux français écrits en macro par les modèles : 1\ière, 3\ième
\newcommand{\ière}{\textsuperscript{re}}\newcommand{\ième}{\textsuperscript{e}}
% \exemple (inventée par les modèles) : étiquette « Exemple : »
\providecommand{\exemple}{\textbf{Exemple~:}\ }
% \square utilisable aussi hors mode math (cases à cocher)
\AtBeginDocument{\let\squareMath\square\renewcommand{\square}{\ensuremath{\squareMath}}}
% Mot ou phrase en anglais à l'intérieur d'un paragraphe français
\newcommand{\eng}[1]{\ifmmode\text{\textit{#1}}\else\begingroup\selectlanguage{english}\itshape #1\endgroup\fi}
\let\esp\eng % alias toléré
\newtcolorbox{exemplebox}[1][]{popbase,colframe=poporange,before upper={\popboxhead{Exemple}{poporange}{#1}}}
\newtcolorbox{definition}[1][]{popbase,colframe=popblue,before upper={\popboxhead{Définition}{popblue}{#1}}}
\newtcolorbox{propriete}[1][]{popbase,colframe=poppurple,before upper={\popboxhead{Propriété}{poppurple}{#1}}}
\newtcolorbox{methode}[1][]{popbase,colframe=popgold,before upper={\popboxhead{Méthode}{popgold}{#1}}}
\newtcolorbox{attention}[1][]{popbase,colframe=poppink,before upper={\popboxhead{Attention}{poppink}{#1}}}
\newtcolorbox{experience}[1][]{popbase,colframe=popblue,colback=popblueL,before upper={\popboxhead{Expérience}{popblue}{#1}}}
% « Le savais-tu ? » : carte violette pleine (contexte, culture, Cameroun)
\newtcolorbox{savaistu}[1][]{popbase,colback=poppurple,colframe=popink,
  fontupper=\popsemi\fontsize{10.4}{14.2}\selectfont\color{white},
  before upper={\poppill{Le savais-tu\,?}{white}{poppurple}\ifx\relax#1\relax\else\hspace{6pt}{\popxbold\color{white}#1}\fi\par\vspace{7pt}}}
% Exercice résolu : énoncé en haut, \tcblower puis la solution sur fond crème
\newcounter{popexo}\newcounter{popexr}
\newtcolorbox{exoresolu}[1][]{popbase,colframe=poporange,colbacklower=popcream,
  segmentation style={popink,line width=1.2pt,dashed},
  before upper={\stepcounter{popexr}\popboxhead{Exercice résolu \thepopexr}{poporange}{#1}},
  before lower={\poppill{Solution}{popink}{popcream}\par\vspace{6pt}}}
% Exercice (non résolu en place) — la correction est dans la partie Corrigés
\newtcolorbox{exercice}[1][]{popbase,colframe=popink,
  before upper={\stepcounter{popexo}\popboxhead{Exercice \thepopexo}{popink}{#1}}}
% Corrigé
\newtcolorbox{corrige}[1][]{popbase,colframe=popgreen,colback=popgreenL,
  before upper={\popboxhead{Corrigé}{popgreen}{#1}}}
% Objectifs / prérequis / bilan
\newtcolorbox{objectifs}{popbase,colframe=popink,colback=poporangeL,
  before upper={\popboxhead{Objectifs de la leçon}{popink}{}}}
\newtcolorbox{prerequis}{popbase,colframe=popink,colback=popblueL,
  before upper={\popboxhead{Prérequis}{popblue}{}}}
\newtcolorbox{bilan}{popbase,colframe=popink,colback=white,boxrule=2.8pt,
  before upper={\popboxhead{Fiche bilan}{poporange}{L'essentiel à connaître pour l'examen}}}
% Figure encadrée avec légende
\newtcolorbox{popfigure}[1][]{enhanced,breakable=false,colback=white,colframe=popline,boxrule=1.4pt,arc=8pt,
  left=10pt,right=10pt,top=10pt,bottom=8pt,before skip=10pt,after skip=12pt,halign=center,
  lower separated=false,halign lower=center,
  fontlower=\popsemi\fontsize{9}{11.5}\selectfont\color{popmuted},
  #1}
\newcommand{\legende}[1]{\par\vspace{6pt}{\popsemi\fontsize{9}{11.5}\selectfont\color{popmuted}#1}}

% ============================================================
% Signature OnBuch+
% ============================================================
\newcommand{\poplogoinline}[1]{{\poplogo\fontsize{#1}{#1}\selectfont\color{popink}OnBuch\color{poporange}+}}
\newcommand{\signatureonbuch}{%
  \begin{tcolorbox}[enhanced,colback=popink,colframe=popink,boxrule=2.2pt,arc=10pt,
      drop shadow=poporange,left=14pt,right=14pt,top=10pt,bottom=10pt,before skip=0pt,after skip=0pt]
    \tikz[baseline=-0.7ex]\node[circle,fill=popgreen,draw=popcream,line width=1.3pt,minimum size=22pt,inner sep=0]{\popcheck{white}};%
    \hspace{9pt}\begin{minipage}[c]{0.62\linewidth}
      {\popxbold\fontsize{12}{13}\selectfont\color{popcream}Signé\ \textbullet\ {\poplogo OnBuch\color{poporange}+}}\par\vspace{2pt}
      {\raggedright\popsemi\fontsize{8}{10.5}\selectfont\color{popline}\MakeUppercase{Équipe pédagogique OnBuch+}\\\MakeUppercase{Conforme au programme officiel MINESEC}\par}
    \end{minipage}\hfill
    {\poplogo\fontsize{22}{22}\selectfont\color{popcream}OnBuch\color{poporange}+}
  \end{tcolorbox}%
}

% ============================================================
% Page de garde
% #1 titre · #2 matière · #3 niveau · #4 module · #5 n° leçon · #6 durée ·
% #7 description / objectifs (texte ou itemize)
% ============================================================
\newcommand{\pagegarde}[7]{%
  \thispagestyle{empty}
  \newgeometry{a4paper,margin=1.7cm,top=1.6cm,bottom=1.4cm}%
  \begin{tikzpicture}[remember picture,overlay]
    \fill[poporange!18] ([xshift=-3.2cm,yshift=-3.6cm]current page.north east) circle (4.6cm);
    \fill[poppurple!14] ([xshift=2.4cm,yshift=7.6cm]current page.south west) circle (3.8cm);
  \end{tikzpicture}%
  {\poplogo\fontsize{30}{30}\selectfont\color{popink}OnBuch\color{poporange}+}\hfill
  \raisebox{0.6ex}{\poppill{Cours signé \textbullet\ Premium}{popink}{popcream}}\par
  \vspace{1pt}\popsquiggle{poppurple}\par
  \vspace{8pt}
  \begin{tcolorbox}[enhanced,colback=poporange,colframe=popink,boxrule=2.8pt,arc=13pt,
      drop shadow=popink,left=18pt,right=18pt,top=12pt,bottom=13pt,after skip=10pt]
    \poppill{#2 \textbullet\ #3}{popink}{popcream}\hspace{5pt}\poppill{#4}{white}{popink}\par\vspace{8pt}
    {\popdisplay\fontsize{14}{14}\selectfont\color{popink}LEÇON #5}\par\vspace{4pt}
    {\raggedright\hyphenpenalty=10000\exhyphenpenalty=10000\popdisplay\fontsize{25}{27}\selectfont\color{popcream}\MakeUppercase{#1}\par}
    \vspace{8pt}
    {\popxbold\fontsize{9.5}{9.5}\selectfont\color{popink}\MakeUppercase{Durée conseillée : #6}}
  \end{tcolorbox}
  \begin{tcolorbox}[enhanced,colback=white,colframe=popink,boxrule=2.2pt,arc=10pt,
      drop shadow=popink,left=16pt,right=16pt,top=10pt,bottom=10pt,after skip=8pt]
    \poppill{Dans cette leçon}{poppurple}{white}\par\vspace{5pt}
    \popsemi\fontsize{9.3}{12.6}\selectfont\color{popink}#7
  \end{tcolorbox}
  \noindent\begin{minipage}[t]{0.487\linewidth}
    \begin{tcolorbox}[enhanced,colback=popgreen,colframe=popink,boxrule=2.2pt,arc=10pt,
        drop shadow=popink,left=11pt,right=11pt,top=10pt,bottom=10pt,height=3.1cm,valign=center]
      \tcbox[colback=white,colframe=popink,boxrule=1.3pt,arc=5pt,boxsep=0pt,nobeforeafter,
        left=3pt,right=3pt,top=3pt,bottom=3pt,valign=center]{\qrcode[height=1.9cm]{https://play.google.com/store/apps/details?id=com.luvvix.onbuch&hl=fr}}%
      \hspace{8pt}\begin{minipage}[c]{0.5\linewidth}
        {\popdisplay\fontsize{11}{12.5}\selectfont\color{white}\MakeUppercase{Télécharge OnBuch}}\par\vspace{4pt}
        {\raggedright\popsemi\fontsize{8.4}{11}\selectfont\color{white}Gratuit \textbullet\ Google Play\\Tuteur IA, annales, concours, résultats\par}
      \end{minipage}
    \end{tcolorbox}
  \end{minipage}\hfill
  \begin{minipage}[t]{0.487\linewidth}
    \begin{tcolorbox}[enhanced,colback=poppurple,colframe=popink,boxrule=2.2pt,arc=10pt,
        drop shadow=popink,left=11pt,right=11pt,top=10pt,bottom=10pt,height=3.1cm,valign=center]
      \tikz[baseline=-0.7ex]\node[circle,fill=white,draw=popink,line width=1.3pt,minimum size=1.6cm,inner sep=0]{\popdisplay\fontsize{24}{24}\selectfont\color{poppurple}+};%
      \hspace{8pt}\begin{minipage}[c]{0.6\linewidth}
        {\popdisplay\fontsize{11}{12.5}\selectfont\color{white}\MakeUppercase{Passe à OnBuch+}}\par\vspace{4pt}
        {\raggedright\popsemi\fontsize{8.4}{11}\selectfont\color{white}Tous les cours signés, Tuteur IA illimité, fiches PDF — onglet OnBuch+ de l'app\par}
      \end{minipage}
    \end{tcolorbox}
  \end{minipage}
  \vspace{8pt}
  \popcontenu
  \vfill
  \signatureonbuch
  \restoregeometry
  \newpage
}

% Rangée « Ce cours contient » (page de garde)
\newcommand{\poptile}[3]{% #1 couleur #2 gros texte #3 légende
  \begin{tcolorbox}[enhanced,colback=#1,colframe=popink,boxrule=2pt,arc=9pt,shadow={2.5pt}{-2.5pt}{0pt}{popink},
      left=6pt,right=6pt,top=9pt,bottom=9pt,halign=center,valign=center,height=1.75cm,width=\linewidth,nobeforeafter]
    {\popdisplay\fontsize{12.5}{13}\selectfont\color{white}\MakeUppercase{#2}}\par\vspace{3pt}
    {\centering\popsemi\fontsize{7.8}{9.5}\selectfont\color{white}#3\par}
  \end{tcolorbox}}
\newcommand{\popcontenu}{%
  \par\noindent{\popxboldls\fontsize{7.5}{7.5}\selectfont\color{popmuted}CE COURS SIGNÉ CONTIENT}\par\vspace{7pt}
  \noindent\begin{minipage}[t]{0.235\linewidth}\poptile{popblue}{Cours}{complet, illustré, pas à pas}\end{minipage}\hfill
  \begin{minipage}[t]{0.235\linewidth}\poptile{poporange}{Méthodes}{et exercices résolus}\end{minipage}\hfill
  \begin{minipage}[t]{0.235\linewidth}\poptile{poppink}{Exercices}{type examen + corrigés}\end{minipage}\hfill
  \begin{minipage}[t]{0.235\linewidth}\poptile{popgreen}{Fiche bilan}{l'essentiel à retenir}\end{minipage}\par}

% Sommaire
\newtcolorbox{sommaire}{enhanced,colback=white,colframe=popink,boxrule=2.2pt,arc=10pt,
  drop shadow=popink,left=18pt,right=18pt,top=13pt,bottom=13pt,before skip=6pt,after skip=20pt,
  before upper={\poppill{Sommaire}{poppurple}{white}\par\vspace{9pt}\popsemi\fontsize{10.6}{14.5}\selectfont\color{popink}}}

% Signature de fin de cours — poussée tout en bas de la dernière page
\newcommand{\findecours}{%
  \par\vfill\nopagebreak
  \begin{center}\popsquiggle{poporange}\end{center}\vspace{4pt}
  \signatureonbuch
}
\AtBeginDocument{\def\llbracket{\mathopen{[\mkern-3.5mu[}}\def\rrbracket{\mathclose{]\mkern-3.5mu]}}} % intervalles d'entiers (glyphe ⟦ absent de la police)
% Glyphes absents de Fira Math : redéfinis à partir de glyphes disponibles
\AtBeginDocument{%
  \def\setminus{\mathbin{\backslash}}\let\smallsetminus\setminus
  \def\vdots{\mathord{\vbox{\baselineskip3.2pt\lineskiplimit0pt\kern4pt\hbox{.}\hbox{.}\hbox{.}}}}%
  \def\ddots{\mathinner{\mkern1mu\raise6.5pt\hbox{.}\mkern2mu\raise3.6pt\hbox{.}\mkern2mu\raise0.7pt\hbox{.}\mkern1mu}}%
  \def\triangle{\mathord{\scalebox{0.9}{$\symup{\Delta}$}}}%
  \def\nabla{\mathord{\raisebox{\depth}{\scalebox{1}[-1]{$\symup{\Delta}$}}}}%
  \def\checkmark{\popcheck{popdark}}%
  \def\star{\mathbin{\ast}}\def\bigstar{\mathord{\ast}}%
  \def\diamond{\mathbin{\scalebox{0.6}{$\Diamond$}}}%
  \def\bigcup{\mathop{\mathchoice{\raisebox{-0.3em}{\scalebox{1.7}{$\cup$}}}{\scalebox{1.25}{$\cup$}}{\cup}{\cup}}\displaylimits}%
  \def\bigcap{\mathop{\mathchoice{\raisebox{-0.3em}{\scalebox{1.7}{$\cap$}}}{\scalebox{1.25}{$\cap$}}{\cap}{\cap}}\displaylimits}%
  \def\lesseqqgtr{\mathrel{\lessgtr}}\def\bigoplus{\mathop{\oplus}\displaylimits}%
}
\providecommand{\ssi}{si et seulement si} % abréviation fréquente
\AtBeginDocument{\providecommand{\wideparen}{\overparen}} % arc de cercle

%% FIGFIX-BEGIN v4 — correctifs globaux des figures (docs/onbuch-generation/tools/figfix)
\usepackage{adjustbox}\usepackage{etoolbox}
% toute figure TikZ/pgfplots plus large que la ligne est réduite à la largeur disponible
\BeforeBeginEnvironment{tikzpicture}{\begin{adjustbox}{max width=\linewidth}}
\AfterEndEnvironment{tikzpicture}{\end{adjustbox}}
% légendes pgfplots lisibles par-dessus les courbes
\pgfplotsset{every axis/.append style={legend style={fill=white,fill opacity=0.92,text opacity=1,draw=black!15,inner sep=3pt}}}
% étiquettes d'arêtes (syntaxe edge["..."]) avec fond blanc
\tikzset{every edge quotes/.append style={fill=white,inner sep=1.2pt}}
%% FIGFIX-END
\begin{document}

\pagegarde{Lesson 5 — Grammar: Review tenses in the previous years + present perfect continuous and future continuous tenses}{Anglais}{Tle A}{Chapter 1 : Family and Social Life}{EN05}{2 h}{Cette leçon de grammaire réactive les temps étudiés depuis la Seconde (present simple, present continuous, past simple, present perfect simple, futurs) avant d'introduire deux temps nouveaux : le present perfect continuous et le future continuous. Elle insiste sur les emplois, les comparaisons avec le français et les erreurs typiques des élèves francophones.
\vspace{4pt}
\begin{multicols}{2}\small
\begin{itemize}
\item À la fin de cette leçon, je sais identifier et conjuguer correctement les principaux temps anglais étudiés les années précédentes
\item À la fin de cette leçon, je sais construire le present perfect continuous (sujet + have/has been + V-ing) à l'affirmative, à la négative et à l'interrogative
\item À la fin de cette leçon, je sais distinguer le present perfect simple du present perfect continuous selon que l'accent porte sur le résultat ou sur la durée
\item À la fin de cette leçon, je sais construire le future continuous (will be + V-ing) et l'utiliser pour une action en cours à un moment du futur
\item À la fin de cette leçon, je sais choisir entre will, be going to et will be + V-ing selon le sens
\item À la fin de cette leçon, je sais repérer les verbes d'état qui ne se mettent pas à la forme continue
\end{itemize}\end{multicols}}

\begin{sommaire}
\begin{multicols}{2}
\begin{enumerate}
\item Observation et réactivation des temps connus
\item Le present perfect continuous : formation
\item Le present perfect continuous : emplois et comparaison avec le present perfect simple
\item Le future continuous : formation et emplois
\item Comparaison français/anglais et pièges des francophones
\item Entraînement guidé et production
\item Activité d'intégration
\item Méthodes et exercices résolus
\item Exercices
\item Corrigés des exercices
\item Fiche bilan
\end{enumerate}
\end{multicols}
\end{sommaire}

\begin{objectifs}
\begin{itemize}
\item À la fin de cette leçon, je sais identifier et conjuguer correctement les principaux temps anglais étudiés les années précédentes
\item À la fin de cette leçon, je sais construire le present perfect continuous (sujet + have/has been + V-ing) à l'affirmative, à la négative et à l'interrogative
\item À la fin de cette leçon, je sais distinguer le present perfect simple du present perfect continuous selon que l'accent porte sur le résultat ou sur la durée
\item À la fin de cette leçon, je sais construire le future continuous (will be + V-ing) et l'utiliser pour une action en cours à un moment du futur
\item À la fin de cette leçon, je sais choisir entre will, be going to et will be + V-ing selon le sens
\item À la fin de cette leçon, je sais repérer les verbes d'état qui ne se mettent pas à la forme continue
\item À la fin de cette leçon, je sais éviter les erreurs fréquentes des francophones (oubli de been, double marqueur, traduction mot à mot)
\item À la fin de cette leçon, je sais employer ces temps dans des phrases et un court paragraphe sur ma vie familiale et scolaire
\end{itemize}
\end{objectifs}

\begin{prerequis}
\begin{itemize}
\item Present simple et present continuous (Seconde)
\item Past simple et past continuous (Première)
\item Present perfect simple avec for, since, just, already, yet (Première)
\item Expression du futur : will et be going to (Première)
\item Formation du participe présent en -ing et ses règles orthographiques
\end{itemize}
\end{prerequis}

\newpage

\coursec{Observation et réactivation des temps connus}

\courssub{Situation d'accroche et problématique}

Imagine : tu rencontres un ami de Bamenda que tu n'as pas vu depuis la classe de Première. Il te dit :

\eng{I have been learning English for six years now, and next year I will be studying at the University of Buea.}

Tu remarques deux formes étranges : \eng{have been learning} et \eng{will be studying}. Pourquoi ne dit-il pas simplement \eng{I learn} ou \eng{I will study} ? Ces deux temps, le \cle{present perfect continuous} et le \cle{future continuous}, permettent d'exprimer la \textbf{durée} et l'\textbf{action en cours dans le futur} — deux notions que le français exprime autrement. Cette leçon t'apprendra à les maîtriser pour le Bac.

Mais avant d'apprendre ces deux nouveaux temps, il faut \textbf{réactiver} tous les temps que tu as étudiés depuis la classe de Seconde. En anglais, chaque temps possède une structure fixe et un emploi précis : c'est la clé de toute la grammaire du Bac.

\courssub{Lecture du texte support : the Ndi family}

Lis attentivement le texte suivant. Pendant la lecture, souligne mentalement toutes les formes verbales conjuguées.

\begin{textebox}[The Ndi family — a text to observe]
The Ndi family lives in Yaoundé. Mr Ndi has worked as a teacher since 2010. Every morning, he drives his children to school before seven o'clock. Right now, his wife is preparing the family meal in the kitchen, and the children are doing their homework quietly.

Last year, they visited their grandparents in Bamenda. They travelled by bus and they spent two wonderful weeks in the hills. Mrs Ndi took many photographs, and the children learned a few words of the local language.

Mr Ndi has already marked forty examination papers this week, but he has not finished yet. He is very tired because he has worked every evening.

Next month, they will travel to Douala for a wedding. The whole family is going to stay at an uncle's house near the market. The children are very excited: they have never seen the sea! Mrs Ndi is going to buy a new dress for the ceremony, and Mr Ndi will probably drive the family car. It will be a long journey, but everybody is looking forward to it.
\end{textebox}

\textbf{Glossaire}

\begin{center}
\begin{tabularx}{\linewidth}{|l|l|X|}\hline
\rowcolor{popblueL} \textbf{Mot anglais} & \textbf{Nature} & \textbf{Traduction} \\ \hline
to mark (papers) & verbe & corriger (des copies) \\ \hline
yet & adverbe & encore (dans une phrase négative) \\ \hline
a wedding & nom & un mariage \\ \hline
the whole family & expression & toute la famille \\ \hline
a journey & nom & un voyage, un trajet \\ \hline
to look forward to & verbe & attendre avec impatience \\ \hline
\end{tabularx}
\end{center}

\textbf{Questions de compréhension (reading)}

\begin{enumerate}
\item Where does the Ndi family live?
\item Since when has Mr Ndi worked as a teacher?
\item What is Mrs Ndi doing right now?
\item Where did the family go last year?
\item What will the family do next month?
\end{enumerate}

\courssub{Repérage et classement des formes verbales}

Reprenons les verbes soulignés dans le texte et classons-les par temps. C'est exactement l'exercice que l'on te demandera de faire mentalement devant n'importe quel texte du Bac.

\begin{center}
\begin{tabularx}{\linewidth}{|X|X|X|}\hline
\rowcolor{popblueL} \textbf{Verbe relevé} & \textbf{Temps} & \textbf{Indice dans le texte} \\ \hline
lives, drives & present simple & habitude : \eng{every morning} \\ \hline
has worked & present perfect & \eng{since 2010} : lien passé-présent \\ \hline
is preparing, are doing & present continuous & \eng{right now} : action en cours \\ \hline
visited, travelled, spent, took, learned & past simple & \eng{last year} : passé accompli \\ \hline
has marked, has not finished, has never seen & present perfect & \eng{already, yet, never} \\ \hline
will travel, will be, will drive & futur avec \eng{will} & \eng{next month} \\ \hline
is going to stay, is going to buy & futur avec \eng{going to} & projet déjà décidé \\ \hline
\end{tabularx}
\end{center}

\begin{propriete}[Un temps = une structure fixe + un emploi précis]
Chaque temps anglais associe une \textbf{structure fixe} (auxiliaire + forme du verbe) à un \textbf{emploi précis}. Le choix du temps dépend du \textbf{rapport du locuteur au temps} :
\begin{itemize}
\item \textbf{habitude} ou vérité générale $\rightarrow$ present simple ;
\item \textbf{action en cours} au moment où l'on parle $\rightarrow$ present continuous ;
\item \textbf{action passée, terminée}, datée $\rightarrow$ past simple ;
\item \textbf{lien entre le passé et le présent} (résultat, expérience, durée avec \eng{since/for}) $\rightarrow$ present perfect ;
\item \textbf{futur} : décision spontanée ou prédiction (\eng{will}) ; projet déjà planifié (\eng{going to}).
\end{itemize}
Contrairement au français, l'anglais ne tolère aucun mélange : on ne peut pas utiliser le present perfect avec une date passée précise (\eng{yesterday}, \eng{last year}, \eng{in 2019}).
\end{propriete}

\begin{methode}[Comment identifier le temps d'un verbe en trois étapes]
\begin{enumerate}
\item \textbf{Repérer l'auxiliaire} : \eng{do/does} (present simple), \eng{be} (continuous), \eng{have/has} (perfect), \eng{will} (futur), \eng{did} ou absence d'auxiliaire (past simple).
\item \textbf{Repérer la forme du verbe principal} : base verbale (live), forme en \eng{-ing} (living), participe passé (lived, worked, seen).
\item \textbf{Croiser avec le marqueur temporel} de la phrase : \eng{every day} (habitude), \eng{now} (en cours), \eng{yesterday} (passé accompli), \eng{since/for/already} (present perfect), \eng{tomorrow/next week} (futur).
\end{enumerate}
Exemple : \eng{has worked} $\rightarrow$ auxiliaire \eng{has} + participe passé $\rightarrow$ present perfect ; le marqueur \eng{since 2010} confirme.
\end{methode}

\courssub{Les marqueurs temporels : les alliés du candidat}

Les \cle{marqueurs temporels} (time markers) sont des adverbes ou expressions qui signalent presque toujours le temps à utiliser. Les mémoriser, c'est gagner des points faciles au Bac.

\begin{center}
\begin{tabularx}{\linewidth}{|l|X|}\hline
\rowcolor{popgreenL} \textbf{Temps} & \textbf{Marqueurs typiques} \\ \hline
Present simple & \eng{every day, always, usually, often, never, on Mondays} \\ \hline
Present continuous & \eng{now, right now, at the moment, today, look!, listen!} \\ \hline
Past simple & \eng{yesterday, last year, ago, in 2015, when I was young} \\ \hline
Present perfect & \eng{since, for, already, just, yet, ever, never, recently} \\ \hline
Futur (\eng{will} / \eng{going to}) & \eng{tomorrow, next month, soon, in 2030, next year} \\ \hline
\end{tabularx}
\end{center}

\begin{exemplebox}[Exemples traduits]
\eng{She has lived here since 2015.} « Elle vit ici depuis 2015. »

\eng{They are playing now.} « Ils jouent en ce moment. »

\eng{He drives his children to school every morning.} « Il conduit ses enfants à l'école tous les matins. »

\eng{We visited Bamenda last year.} « Nous avons visité Bamenda l'année dernière. »

\eng{They will travel to Douala next month.} « Ils voyageront à Douala le mois prochain. »
\end{exemplebox}

\courssub{Tableau récapitulatif des temps déjà étudiés}

\begin{center}
\begin{tabularx}{\linewidth}{|l|X|X|l|}\hline
\rowcolor{popblueL} \textbf{Tense} & \textbf{Structure} & \textbf{Use} & \textbf{Time marker} \\ \hline
Present simple & sujet + base verbale (+ \eng{s} à la 3\textsuperscript{e} pers.) & habitude, vérité générale & \eng{every day} \\ \hline
Present continuous & sujet + \eng{be} + V-\eng{ing} & action en cours maintenant & \eng{now} \\ \hline
Past simple & sujet + V-\eng{ed} ou forme irrégulière & action passée, terminée, datée & \eng{yesterday} \\ \hline
Present perfect & sujet + \eng{have/has} + participe passé & bilan, expérience, durée jusqu'à maintenant & \eng{since, already} \\ \hline
Future (\eng{will}) & sujet + \eng{will} + base verbale & décision spontanée, prédiction & \eng{tomorrow} \\ \hline
Future (\eng{going to}) & sujet + \eng{be going to} + base verbale & projet planifié, signe visible & \eng{next week} \\ \hline
\end{tabularx}
\end{center}

\begin{popfigure}
\begin{tikzpicture}[scale=1]
\draw[-{Stealth[length=3mm]}, thick, popdark] (0,0) -- (13,0);
\foreach \x/\lab in {2/{PASSÉ}, 6.5/{PRÉSENT}, 11/{FUTUR}} {
  \draw[thick, popdark] (\x,-0.15) -- (\x,0.15);
  \node[below, popdark, font=\bfseries\small] at (\x,-0.15) {\lab};
}
\node[above, popblue, align=center, font=\small] at (2,0.25) {past simple\\ \eng{visited}};
\node[above, poporange, align=center, font=\small] at (5.2,0.25) {present perfect\\ \eng{has worked}};
\node[below, popgreen, align=center, font=\small] at (6.5,-1.1) {present simple / continuous\\ \eng{lives / is preparing}};
\node[above, poppurple, align=center, font=\small] at (11,0.25) {\eng{will} / \eng{going to}\\ \eng{will travel}};
\draw[-{Stealth}, thick, poporange] (3.2,0.05) .. controls (4.5,0.9) .. (6.3,0.1);
\node[poporange, font=\scriptsize, align=center] at (4.8,1.15) {lien passé $\rightarrow$ présent};
\end{tikzpicture}
\legende{Frise des temps : le present perfect relie le passé au présent ; le futur se construit avec \eng{will} ou \eng{going to}.}
\end{popfigure}

\begin{attention}[Piège majeur : le présent français n'est pas toujours le présent anglais]
Le présent simple anglais ne traduit pas toujours le présent français. Quand une action \textbf{a commencé dans le passé et continue encore}, le français utilise le présent avec « depuis », mais l'anglais exige le present perfect (ou le present perfect continuous) :

\eng{I have been waiting for two hours.} « J'attends depuis deux heures. »

On ne dit \textbf{jamais} \eng{I am waiting for two hours} : c'est la faute la plus fréquente des francophones au Bac.

Autre piège : \eng{since} + point de départ (\eng{since 2010}) ; \eng{for} + durée (\eng{for two hours}). Ne confonds pas avec « pendant » et « depuis ».
\end{attention}

\begin{exoresolu}[Mini-quiz de conjugaison : mets le verbe au temps correct]
Énoncé — Conjugue le verbe entre parenthèses au temps qui convient.

1. My father \eng{(work)} in Douala. He \eng{(live)} there since 2015.

2. Listen! The students \eng{(sing)} the national anthem.

3. Last week, we \eng{(go)} to the market with our grandmother.

4. I \eng{(already / finish)} my homework, so I can watch television.

5. Tomorrow, the family \eng{(travel)} to Buea for the holidays.

\tcblower
Solution détaillée.

1. \eng{works} : habitude permanente, aucun marqueur de durée $\rightarrow$ present simple (3\textsuperscript{e} personne : \eng{-s}). \eng{has lived} : marqueur \eng{since 2015} $\rightarrow$ present perfect, auxiliaire \eng{has} + participe passé \eng{lived}.

2. \eng{are singing} : marqueur \eng{Listen!} annonce une action en cours $\rightarrow$ present continuous : \eng{be} (\eng{are}) + \eng{sing-ing}.

3. \eng{went} : marqueur \eng{last week} = passé accompli $\rightarrow$ past simple ; \eng{go} est irrégulier : go — went — gone.

4. \eng{have already finished} : marqueur \eng{already} $\rightarrow$ present perfect ; \eng{already} se place entre l'auxiliaire et le participe passé.

5. \eng{will travel} (ou \eng{is going to travel}) : marqueur \eng{tomorrow} $\rightarrow$ futur ; \eng{will} + base verbale.
\end{exoresolu}

\begin{exercice}[À toi de jouer : souligne puis classe]
Dans le texte suivant, souligne les six formes verbales conjuguées, puis classe-les dans un tableau à deux colonnes (verbe / temps) :

\eng{Every Saturday, Amina helps her mother at the market. Today, she is selling tomatoes and onions. She has sold twenty baskets since eight o'clock. Last Saturday, she earned five thousand francs. Next Saturday, she will sell avocados because the season is starting.}

Puis conjugue : 1. They \eng{(play)} football every Friday. 2. We \eng{(not see)} him since January. 3. Look! It \eng{(rain)}. 4. She \eng{(visit)} her uncle two days ago. 5. I \eng{(call)} you tomorrow evening.
\end{exercice}

\begin{corrige}[Exercice « À toi de jouer »]
Verbes du texte : \eng{helps} (present simple), \eng{is selling} (present continuous), \eng{has sold} (present perfect, marqueur \eng{since}), \eng{earned} (past simple, marqueur \eng{last Saturday}), \eng{will sell} (futur avec \eng{will}, marqueur \eng{next Saturday}), \eng{is starting} (present continuous).

Conjugaison : 1. \eng{play} (habitude, \eng{every Friday}). 2. \eng{have not seen} (\eng{since January}). 3. \eng{is raining} (\eng{Look!}). 4. \eng{visited} (\eng{two days ago} = passé accompli). 5. \eng{will call} (\eng{tomorrow}).
\end{corrige}

\begin{aretenir}[L'essentiel de cette section]
\begin{itemize}
\item Chaque temps anglais = une structure fixe + un emploi précis + des marqueurs temporels typiques.
\item Méthode d'identification en trois étapes : auxiliaire $\rightarrow$ forme du verbe $\rightarrow$ marqueur temporel.
\item Le present perfect relie le passé au présent ; il est \textbf{incompatible} avec une date passée précise.
\item « Depuis » + action qui continue = present perfect (continuous), jamais le présent simple anglais.
\end{itemize}
\end{aretenir}

Tu as maintenant réactivé tous les temps connus et tu sais les repérer dans un texte. Dans la prochaine section, nous construirons pas à pas le \cle{present perfect continuous} — la forme \eng{have been learning} de ton ami de Bamenda.

\coursec{Le present perfect continuous : formation}

Dans la section précédente, nous avons réactivé les grands temps étudiés depuis la classe de seconde : le \cle{present simple}, le \cle{present continuous}, le \cle{past simple}, le \cle{past continuous} et le \cle{present perfect simple}. Nous avons notamment rappelé que le present perfect simple (\eng{I have finished}, \eng{she has visited}) relie le passé au présent en insistant sur le \textbf{résultat} d'une action. Mais que faire quand on veut insister non pas sur le résultat, mais sur la \textbf{durée} de l'action, sur le fait qu'elle se poursuit ou vient à peine de s'arrêter ? C'est exactement le rôle du \cle{present perfect continuous}, que nous allons maintenant apprendre à construire rigoureusement.

\courssub{Observation : découvrir la structure dans des phrases authentiques}

\begin{textebox}[At the bus stop in Yaoundé — a short dialogue]
\textbf{Mrs. Ngo Bell:} “Good morning, Paul! You look exhausted. What is happening?”

\textbf{Paul:} “Good morning, Madam. I \textbf{have been waiting} for this bus for forty minutes, and it still has not arrived!”

\textbf{Mrs. Ngo Bell:} “Forty minutes? That is too long. \textbf{Have you been standing} here all this time?”

\textbf{Paul:} “Yes, I have. And look at the sky — it \textbf{has been raining} since early morning, so the bendskins are all full.”

\textbf{Mrs. Ngo Bell:} “I understand. My daughter \textbf{has been learning} English very seriously this year because she is preparing for the Probatoire. She \textbf{has not been sleeping} much lately.”

\textbf{Paul:} “I \textbf{have been revising} for the Bac too. We \textbf{have been working} together with my classmates every evening after school.”

\textbf{Mrs. Ngo Bell:} “That is wonderful. Ah, here comes the bus at last!”

\textbf{Paul:} “Finally! I \textbf{have been dreaming} about this moment for almost an hour!”
\end{textebox}

\textbf{Glossaire :}

\begin{center}
\begin{tabularx}{\linewidth}{|l|l|X|}
\hline
\rowcolor{popblueL} \textbf{Mot anglais} & \textbf{Nature} & \textbf{Traduction} \\
\hline
exhausted & adjectif & épuisé, exténué \\
\hline
to wait & verbe & attendre \\
\hline
bendskin & nom (Cameroun) & moto-taxi \\
\hline
lately & adverbe & ces derniers temps, récemment \\
\hline
to revise & verbe & réviser \\
\hline
at last & locution & enfin \\
\hline
\end{tabularx}
\end{center}

\textbf{Questions d'observation :}

\begin{enumerate}
\item Souligne dans le dialogue tous les verbes au present perfect continuous. Combien y en a-t-il ?
\item Quels sont les \textbf{trois éléments} qui composent chacune de ces formes verbales ?
\item Ces actions sont-elles terminées depuis longtemps, ou sont-elles encore en cours / toutes récentes ?
\item Quels mots indiquent la durée de ces actions ? (\eng{for forty minutes}, \eng{since early morning}, \ldots)
\end{enumerate}

\courssub{De l'observation à la règle : les trois blocs de la structure}

Observons attentivement les phrases clés du dialogue :

\begin{exemplebox}[Phrases d'observation]
\eng{I have been learning English for six years.} — « J'apprends l'anglais depuis six ans. »

\eng{It has been raining since morning.} — « Il pleut depuis ce matin. »

\eng{We have been working together every evening.} — « Nous travaillons ensemble chaque soir. »
\end{exemplebox}

Dans chaque cas, on retrouve exactement \textbf{trois blocs} :

\begin{enumerate}
\item l'auxiliaire \cle{have} conjugué au présent : \eng{have} ou \eng{has} ;
\item le participe passé du verbe \eng{be}, c'est-à-dire \cle{been} ;
\item le \cle{participe présent} du verbe principal, c'est-à-dire le verbe en \textbf{-ing}.
\end{enumerate}

\begin{propriete}[Formation du present perfect continuous — forme affirmative]
Le \cle{present perfect continuous} (aussi appelé \emph{present perfect progressive}) se forme ainsi :

\begin{center}
\textbf{Sujet + have / has + been + verbe en -ing}
\end{center}

\begin{itemize}
\item \eng{have} s'emploie avec \eng{I}, \eng{you}, \eng{we}, \eng{they} ;
\item \eng{has} s'emploie avec \eng{he}, \eng{she}, \eng{it} (troisième personne du singulier) ;
\item \eng{been} est \textbf{invariable} : il ne change jamais ;
\item le verbe principal prend toujours la terminaison \textbf{-ing}, sans exception de conjugaison.
\end{itemize}

Ce temps relie une action \textbf{commencée dans le passé} au \textbf{moment présent}, en insistant sur sa \textbf{durée} ou sa \textbf{continuité}.
\end{propriete}

\begin{popfigure}
\begin{center}
\begin{tikzpicture}[
  bloc/.style={rectangle, rounded corners=3pt, draw=popline, thick, align=center, minimum height=1.1cm, inner sep=6pt},
  fleche/.style={-{Stealth[length=2.5mm]}, thick, popmuted},
  etiquette/.style={align=center, font=\small, text=popmuted}
]
\node[bloc, fill=popblueL, text=popdark, align=center] (s) at (0,0){\textbf{I / She}\\ \textbf{We / They}};
\node[bloc, fill=poporangeL, text=popdark] (h) at (3.4,0) {\textbf{have / has}};
\node[bloc, fill=popgreenL, text=popdark] (b) at (6.6,0) {\textbf{been}};
\node[bloc, fill=poppinkL, text=popdark, align=center] (v) at (9.8,0){\textbf{learn-ing}\\ \textbf{rain-ing}};

\draw[fleche] (s) -- (h);
\draw[fleche] (h) -- (b);
\draw[fleche] (b) -- (v);

\node[etiquette] at (0,-1.15) {sujet};
\node[etiquette, align=center] at (3.4,-1.15){auxiliaire \emph{have}\\ conjugué au présent};
\node[etiquette, align=center] at (6.6,-1.15){participe passé\\ de \emph{be}};
\node[etiquette, align=center] at (9.8,-1.15){participe présent\\ du verbe (-ing)};
\end{tikzpicture}
\end{center}
\legende{Schéma de structure du present perfect continuous : trois blocs obligatoires, dans cet ordre.}
\end{popfigure}

\courssub{Tableau de conjugaison complet}

\begin{popfigure}
\begin{center}
\begin{tabular}{|l|l|l|}
\hline
\rowcolor{popblueL} \textbf{Personne} & \textbf{Forme affirmative} & \textbf{Traduction} \\
\hline
I & have been working & je travaille (depuis\ldots) \\
\hline
You & have been working & tu travailles (depuis\ldots) \\
\hline
He / She / It & \textbf{has} been working & il / elle travaille (depuis\ldots) \\
\hline
We & have been working & nous travaillons (depuis\ldots) \\
\hline
You & have been working & vous travaillez (depuis\ldots) \\
\hline
They & have been working & ils / elles travaillent (depuis\ldots) \\
\hline
\end{tabular}
\end{center}
\legende{Mini-tableau de conjugaison : seul l'auxiliaire change (have / has) ; been + V-ing restent identiques partout.}
\end{popfigure}

\begin{aretenir}[Une seule chose à conjuguer !]
Bonne nouvelle : au present perfect continuous, \textbf{seul l'auxiliaire} \eng{have} se conjugue (\eng{have} / \eng{has}). Le reste — \eng{been} + verbe en -ing — est \textbf{identique pour toutes les personnes}. Il n'y a donc qu'un seul choix à faire : \eng{have} ou \eng{has}.
\end{aretenir}

\courssub{La forme négative}

\begin{propriete}[Forme négative]
On place \cle{not} juste après l'auxiliaire \eng{have} / \eng{has} :

\begin{center}
\textbf{Sujet + have / has + not + been + verbe en -ing}
\end{center}

Formes contractées courantes : \eng{haven't been working}, \eng{hasn't been working}.

\eng{She has not been sleeping well lately.} — « Elle ne dort pas bien ces derniers temps. »

\eng{They haven't been listening to the teacher.} — « Ils n'écoutent pas le professeur (depuis un moment). »
\end{propriete}

\begin{exemplebox}[Exemples traduits — forme négative]
\eng{I have not been feeling well since Monday.} — « Je ne me sens pas bien depuis lundi. »

\eng{He hasn't been attending classes regularly.} — « Il ne suit pas les cours régulièrement (ces derniers temps). »

\eng{We haven't been wasting our time.} — « Nous n'avons pas été en train de perdre notre temps. »
\end{exemplebox}

\courssub{La forme interrogative}

\begin{propriete}[Forme interrogative]
On inverse l'auxiliaire \eng{have} / \eng{has} et le sujet ; \eng{been} + V-ing restent après le sujet :

\begin{center}
\textbf{Have / Has + sujet + been + verbe en -ing ?}
\end{center}

\eng{Have you been waiting long?} — « Tu attends depuis longtemps ? »

\eng{Has she been studying for the exam?} — « Est-ce qu'elle révise pour l'examen ? »

Avec un mot interrogatif, celui-ci se place en tête : \eng{How long have you been waiting?} — « Depuis combien de temps attends-tu ? »
\end{propriete}

\begin{methode}[Former une question au present perfect continuous en 4 étapes]
Pour transformer \eng{You have been learning English for six years} en question :

\begin{enumerate}
\item \textbf{Repère} l'auxiliaire : c'est \eng{have}.
\item \textbf{Inverse} l'auxiliaire et le sujet : \eng{have you}.
\item \textbf{Garde} \eng{been} + V-ing sans rien changer : \eng{been learning}.
\item \textbf{Ajoute} le mot interrogatif au début si nécessaire : \eng{How long have you been learning English?}
\end{enumerate}

Résultat : \eng{How long have you been learning English?} — « Depuis combien de temps apprends-tu l'anglais ? »
\end{methode}

\courssub{Rappel orthographique : comment former le participe présent en -ing}

Le troisième bloc de la structure exige le verbe en \textbf{-ing}. Rappelons les règles d'orthographe vues dans les classes précédentes :

\begin{center}
\begin{tabularx}{\linewidth}{|l|l|X|}
\hline
\rowcolor{popblueL} \textbf{Règle} & \textbf{Exemples} & \textbf{Explication} \\
\hline
Cas général : on ajoute -ing & work $\to$ working, rain $\to$ raining, study $\to$ studying & Aucun changement. Attention : le \textbf{y} ne change jamais devant -ing. \\
\hline
Verbe en -e muet : on supprime le e & make $\to$ making, write $\to$ writing, have $\to$ having & Le -e final muet tombe devant -ing. \\
\hline
CVC : on double la consonne finale & run $\to$ running, sit $\to$ sitting, swim $\to$ swimming & Verbe d'une syllabe : consonne + voyelle + consonne. \\
\hline
Verbe en -ie & lie $\to$ lying, die $\to$ dying & -ie devient -y devant -ing. \\
\hline
\end{tabularx}
\end{center}

\begin{attention}[Le y ne change jamais devant -ing !]
Contrairement au present simple (\eng{he studies}), le \textbf{y} reste intact devant -ing : \eng{study $\to$ studying} (et non \emph{studing}), \eng{cry $\to$ crying}, \eng{try $\to$ trying}. Mais attention au doublement : \eng{run $\to$ running} avec \textbf{deux n}.
\end{attention}

\courssub{Les contractions à l'oral}

À l'oral et dans l'écrit informel, l'auxiliaire \eng{have} / \eng{has} se contracte presque toujours avec le sujet :

\begin{center}
\begin{tabularx}{\linewidth}{|X|X|X|}
\hline
\rowcolor{popblueL} \textbf{Forme pleine} & \textbf{Forme contractée} & \textbf{Exemple} \\
\hline
I have been & I've been & \eng{I've been waiting for you.} \\
\hline
You have been & You've been & \eng{You've been working hard.} \\
\hline
He / She has been & He's / She's been & \eng{She's been crying.} \\
\hline
We have been & We've been & \eng{We've been revising all week.} \\
\hline
They have been & They've been & \eng{They've been playing football.} \\
\hline
\end{tabularx}
\end{center}

Prononciation : \eng{I've been} se dit « aïv bine », \eng{she's been} se dit « chiz bine ». À l'écrit formel (rédaction, lettre officielle), on préfère les formes pleines.

\courssub{Comparaison avec le français}

\begin{center}
\begin{tabularx}{\linewidth}{|X|X|X|}
\hline
\rowcolor{popblueL} \textbf{Français} & \textbf{English} & \textbf{Remarque} \\
\hline
J'apprends l'anglais depuis six ans. & \eng{I have been learning English for six years.} & Le français utilise le \textbf{présent} + « depuis » ; l'anglais exige le present perfect continuous. \\
\hline
Il pleut depuis ce matin. & \eng{It has been raining since morning.} & Jamais de présent simple anglais avec \eng{since}. \\
\hline
Nous révisons pour le Bac depuis toute la semaine. & \eng{We have been revising for the Bac all week.} & « Depuis + durée » = \eng{for} ; « depuis + point de départ » = \eng{since}. \\
\hline
Tu as les yeux rouges — tu as pleuré ? & \eng{Your eyes are red — have you been crying?} & L'anglais insiste sur l'action qui explique le résultat visible. \\
\hline
\end{tabularx}
\end{center}

\begin{exemplebox}[Exemples traduits supplémentaires]
\eng{We have been revising for the Bac all week.} — « Nous révisons pour le Bac depuis toute la semaine. »

\eng{She has not been sleeping well lately.} — « Elle ne dort pas bien ces derniers temps. »

\eng{The students have been preparing the debate since last month.} — « Les élèves préparent le débat depuis le mois dernier. »

\eng{My father has been working in Douala for twenty years.} — « Mon père travaille à Douala depuis vingt ans. »
\end{exemplebox}

\courssub{Exercice résolu et pièges à éviter}

\begin{exoresolu}[Mets les verbes au present perfect continuous]
\textbf{Énoncé.} Conjugue les verbes entre parenthèses au present perfect continuous.

1. I \ldots\ldots\ldots (study) English for six years.

2. She \ldots\ldots\ldots (not / sleep) well lately.

3. How long \ldots\ldots\ldots you \ldots\ldots\ldots (wait) for the bus?

4. They \ldots\ldots\ldots (run) for twenty minutes; they are tired.

5. It \ldots\ldots\ldots (rain) since this morning.

\tcblower
\textbf{Solution détaillée.}

1. \eng{I \textbf{have been studying} English for six years.} — Sujet \eng{I} $\to$ \eng{have} ; \eng{study} garde son \eng{y} devant -ing : \eng{studying}.

2. \eng{She \textbf{has not been sleeping} well lately.} — Sujet \eng{she} $\to$ \eng{has} ; forme négative : \eng{not} après l'auxiliaire.

3. \eng{How long \textbf{have} you \textbf{been waiting} for the bus?} — Inversion de \eng{have} et du sujet \eng{you} ; \eng{been waiting} inchangé.

4. \eng{They \textbf{have been running} for twenty minutes.} — \eng{run} double sa consonne finale (CVC) : \eng{running}.

5. \eng{It \textbf{has been raining} since this morning.} — Sujet \eng{it} $\to$ \eng{has}.
\end{exoresolu}

\begin{attention}[L'erreur n°~1 des francophones : l'oubli de \emph{been}]
Ne jamais oublier \cle{been} entre l'auxiliaire et le verbe en -ing :

\begin{itemize}
\item FAUX : \eng{I have studying English for six years.}
\item CORRECT : \eng{I have \textbf{been} studying English for six years.}
\end{itemize}

Autres pièges fréquents :

\begin{itemize}
\item Ne confonds pas \eng{been} (participe passé de \eng{be}) et \eng{being} : on dit \eng{I have been working}, jamais \emph{I have being working}.
\item N'oublie pas le \textbf{-ing} : \eng{she has been work} est faux ; il faut \eng{she has been working}.
\item Ne traduis pas « depuis » par un présent simple : \emph{I learn English since six years} est une faute grave ; dis \eng{I have been learning English for six years}.
\end{itemize}
\end{attention}

\begin{savaistu}[Un temps très vivant à l'oral]
En anglais parlé, britannique comme américain, le present perfect continuous est extrêmement fréquent pour \textbf{expliquer une situation visible} au moment où l'on parle. Quelques exemples typiques :

\begin{itemize}
\item \eng{Your eyes are red — have you been crying?} — « Tu as les yeux rouges — tu as pleuré ? »
\item \eng{You are all wet! Have you been walking in the rain?} — « Tu es trempé ! Tu as marché sous la pluie ? »
\item \eng{Why are you so tired? — I have been playing football all afternoon.} — « Pourquoi es-tu si fatigué ? — J'ai joué au football tout l'après-midi. »
\end{itemize}

C'est le temps de la \textbf{déduction} : on observe un résultat présent et on devine l'action qui vient de se dérouler. Tu l'entendras constamment dans les séries et les chansons anglophones !
\end{savaistu}

\begin{exercice}[À toi de jouer]
1. Conjugue au present perfect continuous : (a) We \ldots (wait) for the results since Monday. (b) He \ldots (not / attend) classes regularly. (c) \ldots you \ldots (revise) for the Probatoire?

2. Mets à la forme interrogative : \eng{She has been working in Buea for five years.} (avec \eng{How long})

3. Corrige les erreurs : (a) \emph{I have being learning French.} (b) \emph{They has been playing since morning.} (c) \emph{She have been study all night.}
\end{exercice}

\begin{corrige}[Exercice « À toi de jouer »]
1. (a) \eng{We \textbf{have been waiting} for the results since Monday.} (b) \eng{He \textbf{has not been attending} classes regularly.} (c) \eng{\textbf{Have} you \textbf{been revising} for the Probatoire?}

2. \eng{How long \textbf{has she been working} in Buea?} — On inverse \eng{has} et le sujet \eng{she}, on garde \eng{been working}.

3. (a) \eng{I have \textbf{been} learning French.} — participe passé de \eng{be} = \eng{been}, pas \eng{being}. (b) \eng{They \textbf{have} been playing since morning.} — \eng{they} exige \eng{have}. (c) \eng{She \textbf{has been studying} all night.} — \eng{she} exige \eng{has}, et le verbe principal doit être en -ing.
\end{corrige}

Maintenant que la \textbf{formation} du present perfect continuous est parfaitement maîtrisée — trois blocs, un seul choix (\eng{have} ou \eng{has}), et le fameux \eng{been} à ne jamais oublier —, nous pouvons passer à la section suivante : les \textbf{emplois} précis de ce temps et sa comparaison avec le present perfect simple, dont la nuance est un classique des épreuves du Baccalauréat.

\coursec{Le present perfect continuous : emplois et comparaison avec le present perfect simple}

Dans la section précédente, nous avons appris à \emph{former} le present perfect continuous : \textbf{have/has + been + verbe en -ing} (\eng{I have been working}, \eng{she has been studying}). Savoir conjuguer un temps ne suffit pas : il faut maintenant comprendre \emph{quand} l'employer, et surtout comment le distinguer de son proche voisin, le \emph{present perfect simple}. C'est précisément l'objet de cette section, la plus importante de la leçon.

\courssub{Observation : un texte pour découvrir les emplois}

\begin{textebox}[A tired family — reading text]
It is Saturday evening in the Ngoa-Ekelle neighbourhood of Yaoundé. Mrs Biya looks at her husband and smiles : “You look exhausted ! What have you been doing all day ?” Mr Biya sits down heavily. “I have been repairing the roof since eight o'clock this morning. It has been raining a lot recently, and the water was coming in.” Their daughter Sandrine enters the room with red eyes. “Have you been crying ?” her mother asks. “No, I have been studying for my English test. I have been reading this grammar book for three hours !” Her little brother Alain laughs : “I have finished my homework. I have written two essays and I have cleaned my room.” Mrs Biya serves dinner. “Well, I have been cooking since four o'clock, so I hope everyone is hungry !” Outside, the rain has stopped, but the family has been listening to it all afternoon.
\end{textebox}

\begin{definition}[Glossaire]
\begin{itemize}
\item \cle{exhausted} (adj.) : épuisé, très fatigué
\item \cle{roof} (n.) : le toit
\item \cle{to come in} (phr. v.) : entrer, s'infiltrer (pour l'eau)
\item \cle{red eyes} : les yeux rouges
\item \cle{hungry} (adj.) : affamé
\end{itemize}
\end{definition}

\textbf{Questions de compréhension :}
\begin{enumerate}
\item What has Mr Biya been doing since eight o'clock ?
\item Why are Sandrine's eyes red ?
\item What has Alain finished ? How many essays has he written ?
\item How long has Mrs Biya been cooking ?
\end{enumerate}

\begin{corrige}[Questions de compréhension]
1. He has been repairing the roof. 2. Because she has been studying (and reading her grammar book for three hours). 3. He has finished his homework ; he has written two essays. 4. She has been cooking since four o'clock.
\end{corrige}

Observons maintenant les phrases du texte. Certaines décrivent des actions \emph{en cours ou dont la durée compte} (\eng{I have been repairing the roof since eight o'clock}) ; d'autres annoncent un \emph{résultat accompli} (\eng{I have written two essays}). C'est toute la différence entre les deux temps.

\courssub{Les trois grands emplois du present perfect continuous}

\textbf{Emploi 1 : une action commencée dans le passé qui continue encore (avec \eng{for} et \eng{since})}

C'est l'emploi le plus fréquent. L'action a commencé dans le passé et \emph{n'est pas terminée} au moment où l'on parle. On précise la durée avec \cle{for} (durée : \eng{for two hours}, \eng{for ten years}) ou le point de départ avec \cle{since} (\eng{since Monday}, \eng{since 2015}, \eng{since I was a child}).

\begin{exemplebox}[Action qui continue encore]
\eng{I have been learning English for six years.} « J'apprends l'anglais depuis six ans. » (j'apprends encore)\par
\eng{They have been living in Douala since 2018.} « Ils habitent à Douala depuis 2018. »\par
\eng{She has been waiting for a bendskin for twenty minutes.} « Elle attend une moto-taxi depuis vingt minutes. »\par
\eng{We have been working on this project all week.} « Nous travaillons sur ce projet toute la semaine. »
\end{exemplebox}

\textbf{Emploi 2 : une action récente dont le résultat est visible maintenant}

L'action vient de se terminer (ou est presque terminée), et on en voit les \emph{traces physiques} au moment présent : fatigue, vêtements mouillés, yeux rouges, mains sales.

\begin{exemplebox}[Résultat visible maintenant]
\eng{“You are wet !” — “Yes, I have been walking in the rain.”} « Tu es trempé ! — Oui, j'ai marché sous la pluie. »\par
\eng{His hands are dirty because he has been repairing his motorbike.} « Il a les mains sales parce qu'il a réparé sa moto. »\par
\eng{I am tired ; I have been playing football all afternoon.} « Je suis fatigué ; j'ai joué au football tout l'après-midi. »\par
\eng{Why are your eyes red ? — I have been crying.} « Pourquoi as-tu les yeux rouges ? — J'ai pleuré. »
\end{exemplebox}

\textbf{Emploi 3 : une action répétée sur une période récente}

Le temps peut aussi exprimer une action \emph{répétée, habituelle} sur une période récente qui continue jusqu'au présent, souvent avec \eng{all day}, \eng{all week}, \eng{lately}, \eng{recently}.

\begin{exemplebox}[Action répétée]
\eng{He has been calling her all day.} « Il l'a appelée toute la journée. » (plusieurs appels)\par
\eng{It has been raining a lot recently.} « Il a beaucoup plu récemment. »\par
\eng{She has been visiting her grandmother every weekend lately.} « Elle rend visite à sa grand-mère chaque week-end ces derniers temps. »
\end{exemplebox}

\begin{propriete}[Règle fondamentale : durée contre résultat]
Le \cle{present perfect continuous} (\eng{have/has been + V-ing}) met l'accent sur la \textbf{durée}, le \textbf{processus} ou l'action \textbf{non terminée}. Le \cle{present perfect simple} (\eng{have/has + participe passé}) met l'accent sur le \textbf{résultat}, la \textbf{quantité accomplie} ou l'action \textbf{terminée}. Les deux temps relient le passé au présent, mais l'angle de vue change : \emph{comment l'action s'est déroulée} (continuous) contre \emph{ce qui a été accompli} (simple).
\end{propriete}

\courssub{Comparaison systématique : simple contre continuous}

Comparons les deux temps sur des paires de phrases identiques :

\begin{exemplebox}[Paires contrastives]
\eng{I have been reading this novel for a month.} « Je lis ce roman depuis un mois. » (\emph{pas fini} — je suis encore en train de le lire)\par
\eng{I have read this novel.} « J'ai lu ce roman. » (\emph{fini} — je connais la fin)\par
\eng{I have been writing letters all morning.} « J'écris des lettres depuis ce matin. » (durée, activité en cours)\par
\eng{I have written three letters.} « J'ai écrit trois lettres. » (résultat chiffré, tâche accomplie)\par
\eng{Someone has been eating my mangoes !} « Quelqu'un a mangé de mes mangues ! » (il en reste peut-être, mais je vois les traces)\par
\eng{Someone has eaten all my mangoes !} « Quelqu'un a mangé toutes mes mangues ! » (résultat : il n'en reste plus)
\end{exemplebox}

\begin{center}
\begin{tabularx}{\linewidth}{|X|X|}
\hline
\rowcolor{popblueL} \textbf{Present perfect simple (résultat)} & \textbf{Present perfect continuous (durée)} \\
\hline
\eng{I have read the book.} (fini) & \eng{I have been reading the book.} (pas fini) \\
\hline
\eng{She has written three letters.} (quantité) & \eng{She has been writing letters all morning.} (durée) \\
\hline
\eng{He has painted the wall.} (le mur est prêt) & \eng{He has been painting the wall.} (il est fatigué, couvert de peinture) \\
\hline
\eng{We have visited Bamenda twice.} (nombre de fois) & \eng{We have been visiting our cousins in Bamenda.} (activité récente) \\
\hline
\eng{How many exercises have you done ?} & \eng{How long have you been doing your homework ?} \\
\hline
\end{tabularx}
\end{center}

\begin{popfigure}
\begin{tikzpicture}[mindmap, grow cyclic, every node/.style={concept, align=center, font=\small},
root concept/.append style={concept color=popblue, font=\small\bfseries},
level 1/.append style={level distance=3.6cm, sibling angle=120},
level 1 concept/.append style={concept color=poporange}]
\node[root concept, align=center]{Present perfect\\continuous};
\node[level 1 concept, align=center]{Duration\\for / since\\all day}
  child[grow=150, concept color=popgreen]{node[concept, align=center, font=\scriptsize]{I have been\\waiting for\\an hour.}};
\node[level 1 concept, align=center]{Visible result\\now}
  child[grow=30, concept color=popgreen]{node[concept, align=center, font=\scriptsize]{You are wet :\\you have been\\walking in\\the rain.}};
\node[level 1 concept, align=center]{Repeated action\\lately / recently}
  child[grow=-90, concept color=popgreen]{node[concept, align=center, font=\scriptsize]{He has been\\calling her\\all day.}};
\end{tikzpicture}
\legende{Carte mentale : les trois emplois du present perfect continuous.}
\end{popfigure}

\courssub{L'exception majeure : les verbes d'état}

\begin{attention}[Les verbes d'état refusent la forme continue]
Les verbes qui expriment un \emph{état} (et non une action) — \cle{know}, \cle{believe}, \cle{like}, \cle{love}, \cle{want}, \cle{need}, \cle{understand}, \cle{own}, \cle{belong}, \cle{seem}, \cle{hate}, \cle{remember} — ne se mettent \textbf{jamais} à la forme continue. Avec \eng{for} et \eng{since}, on utilise donc obligatoirement le present perfect \emph{simple} :\par
\eng{I have known him since childhood.} « Je le connais depuis l'enfance. » (jamais \emph{I have been knowing him})\par
\eng{She has wanted a bicycle for years.} « Elle veut un vélo depuis des années. » (jamais \emph{she has been wanting})\par
\eng{We have understood the lesson.} « Nous avons compris la leçon. »
\end{attention}

\textbf{Cas particulier : \eng{live} et \eng{work}.} Ces deux verbes acceptent les deux temps sans grande différence de sens : \eng{I have lived in Buea for ten years} et \eng{I have been living in Buea for ten years} sont tous deux corrects ; le continuous insiste simplement un peu plus sur le caractère temporaire ou en cours.

\courssub{Méthode, comparaison avec le français et entraînement}

\begin{methode}[Comment choisir entre simple et continuous]
\begin{enumerate}
\item Cherche dans la phrase une \textbf{quantité} ou un \textbf{résultat chiffré} (\eng{three letters}, \eng{twice}, \eng{all my mangoes}, \eng{the book} = fini) : utilise le \textbf{present perfect simple}.
\item Cherche une expression de \textbf{durée} ou d'activité en cours (\eng{all day}, \eng{for hours}, \eng{since this morning}, un résultat physique visible) : utilise le \textbf{present perfect continuous}.
\item Vérifie le verbe : s'il s'agit d'un \textbf{verbe d'état} (\eng{know}, \eng{like}, \eng{want}, \eng{believe}, \eng{understand}, \eng{own}), utilise toujours le \textbf{simple}, même avec \eng{for} ou \eng{since}.
\item Cas libre : avec \eng{live} et \eng{work}, les deux sont possibles.
\end{enumerate}
\end{methode}

\begin{center}
\begin{tabularx}{\linewidth}{|X|X|X|}
\hline
\rowcolor{popblueL} Français & English & Remarque \\
\hline
J'attends depuis une heure. & I have been waiting for an hour. & Le français utilise le présent + « depuis » ; l'anglais exige le present perfect. \\
\hline
Il pleut depuis ce matin. & It has been raining since this morning. & Même remarque. \\
\hline
Je le connais depuis dix ans. & I have known him for ten years. & Verbe d'état : simple obligatoire. \\
\hline
J'ai écrit trois lettres. & I have written three letters. & Résultat chiffré : simple. \\
\hline
\end{tabularx}
\end{center}

\begin{savaistu}[Pourquoi les francophones disent « I am waiting for an hour »]
En français, on dit naturellement « J'attends \emph{depuis} une heure » avec un simple présent. Beaucoup d'élèves traduisent donc mot à mot : \emph{I am waiting for an hour} — phrase incorrecte en anglais ! L'anglais refuse le présent avec \eng{for} ou \eng{since} : il faut impérativement un present perfect, simple ou continuous. Retiens l'équation : \textbf{présent français + « depuis » = present perfect anglais}.
\end{savaistu}

\begin{exoresolu}[Choisis le bon temps]
Complète avec le present perfect simple ou continuous du verbe entre parenthèses :\par
1. I \dots (study) for this exam for three weeks, and I am still not ready.\par
2. She \dots (write) five applications this morning. Look at them on the table !\par
3. We \dots (know) the Mbappe family since we were children.\par
4. Look at your clothes ! You \dots (play) in the mud again !\par
5. He \dots (work) in this bank for twenty years. He loves his job.
\tcblower
\textbf{Solution détaillée :}\par
1. \eng{have been studying} — durée (\eng{for three weeks}) et action non terminée (\eng{still not ready}) : continuous.\par
2. \eng{has written} — quantité accomplie (\eng{five applications}) : simple.\par
3. \eng{have known} — verbe d'état : simple obligatoire, jamais \emph{have been knowing}.\par
4. \eng{have been playing} — résultat visible maintenant (vêtements sales) : continuous.\par
5. \eng{has worked} ou \eng{has been working} — \eng{work} accepte les deux temps sans différence de sens réelle.
\end{exoresolu}

\begin{attention}[Erreurs fréquentes des francophones]
\begin{itemize}
\item \emph{I am learning English since six years.} Faux ! Avec \eng{since/for}, jamais de présent : \eng{I have been learning English for six years.} (et \eng{for six years}, car \eng{since} s'emploie avec un point de départ : \eng{since 2019}).
\item \emph{I have been knowing her for years.} Faux ! Verbe d'état : \eng{I have known her for years.}
\item \emph{I have been writing three letters.} Maladroit : la quantité accomplie appelle le simple : \eng{I have written three letters.}
\item Oublier \eng{been} : \emph{I have waiting for an hour} est impossible ; la forme exige \eng{have been + V-ing}.
\end{itemize}
\end{attention}

\begin{exercice}[À toi de jouer]
Mets les verbes entre parenthèses au present perfect simple ou continuous, puis traduis chaque phrase en français :\par
1. They \dots (build) their house for two years ; it is almost finished.\par
2. I \dots (lose) my keys ! I cannot enter the house.\par
3. She \dots (believe) in God since her childhood.\par
4. The students \dots (make) too much noise all morning ; the headmaster is angry.\par
5. How many goals \dots Eto'o \dots (score) this season ?
\end{exercice}

\begin{corrige}[À toi de jouer]
1. \eng{have been building} — « Ils construisent leur maison depuis deux ans ; elle est presque finie. » (durée, action non terminée). 2. \eng{have lost} — « J'ai perdu mes clés ! » (résultat présent). 3. \eng{has believed} — « Elle croit en Dieu depuis son enfance. » (verbe d'état). 4. \eng{have been making} — « Les élèves font trop de bruit depuis ce matin. » (durée + résultat visible : le proviseur est en colère). 5. \eng{has Eto'o scored} — « Combien de buts Eto'o a-t-il marqués cette saison ? » (quantité : simple).
\end{corrige}

\begin{aretenir}[L'essentiel de la section]
Le present perfect continuous insiste sur la \textbf{durée}, le \textbf{processus} ou la \textbf{répétition} ; le present perfect simple insiste sur le \textbf{résultat} ou la \textbf{quantité}. Les \textbf{verbes d'état} (\eng{know}, \eng{like}, \eng{want}, \eng{believe}, \eng{understand}, \eng{own}) n'admettent que le simple. \eng{Live} et \eng{work} acceptent les deux. Enfin, « depuis + présent » en français se traduit toujours par un present perfect en anglais.
\end{aretenir}

\coursec{Le future continuous : formation et emplois}

Dans la section précédente, nous avons étudié le \cle{present perfect continuous} (\eng{have/has been + V-ing}), qui décrit une action commencée dans le passé et encore en cours au moment présent. Nous restons dans la famille des temps « continus » — ceux qui mettent l'accent sur le \emph{déroulement} d'une action — mais nous nous projetons cette fois dans l'avenir : bienvenue au \cle{future continuous}, le temps de l'action \emph{en train de se dérouler} à un moment précis du futur.

\courssub{Observation : des exemples pour découvrir la règle}

\begin{exemplebox}[Observez et comparez]
Lisez attentivement les phrases suivantes et repérez ce qui se répète :

\begin{enumerate}
\item \eng{This time tomorrow, I will be sitting in the exam room.} « Demain à la même heure, je serai en train de composer dans la salle d'examen. »
\item \eng{Next year, she will be studying at university.} « L'année prochaine, elle sera en train d'étudier à l'université. »
\item \eng{At 8 p.m., we will be having dinner.} « À 20 heures, nous serons en train de dîner. »
\item \eng{Will you be using your phone tonight?} « Est-ce que tu utiliseras ton téléphone ce soir ? »
\end{enumerate}

\textbf{Questions d'observation :}
\begin{itemize}
\item Quel mot revient avant chaque verbe ? (\eng{will be})
\item Quelle terminaison porte le verbe principal ? (\eng{-ing})
\item Le mot \eng{will} change-t-il selon le sujet (\eng{I}, \eng{she}, \eng{we}, \eng{you}) ? (Non : il est \emph{invariable}.)
\end{itemize}
\end{exemplebox}

On remarque immédiatement la structure : \eng{will be + verbe en -ing}. Le sens est également clair : chaque phrase décrit une action qui sera \emph{en cours}, « en train de » se faire, à un moment donné du futur. Le français exprime souvent cette idée avec « être en train de » au futur, ou simplement avec le futur simple.

\courssub{Formation du future continuous}

\begin{propriete}[Le future continuous : forme et sens]
Le \cle{future continuous} (ou \emph{future progressive}) se forme ainsi :

\begin{center}
\textbf{Sujet + will be + verbe en -ing}
\end{center}

\begin{itemize}
\item \eng{will} est \textbf{invariable} : il ne change jamais, quelle que soit la personne (\eng{I will be}, \eng{he will be}, \eng{they will be}).
\item \eng{be} reste toujours à la base verbale : jamais \eng{will is} ni \eng{will are}.
\item Le verbe principal prend toujours la terminaison \eng{-ing} (règles habituelles : \eng{work $\to$ working}, \eng{write $\to$ writing}, \eng{sit $\to$ sitting}).
\item Il décrit une \textbf{action qui sera en train de se dérouler} à un moment donné du futur.
\end{itemize}
\end{propriete}

\textbf{Les trois formes : affirmative, négative, interrogative}

\begin{center}
\begin{tabular}{|l|l|l|}
\hline
\rowcolor{popblueL} \textbf{Forme} & \textbf{Structure} & \textbf{Exemple} \\
\hline
Affirmative & Sujet + will be + V-ing & \eng{She will be working.} \\
\hline
Négative & Sujet + will not (won't) be + V-ing & \eng{She won't be working.} \\
\hline
Interrogative & Will + sujet + be + V-ing ? & \eng{Will she be working?} \\
\hline
\end{tabular}
\end{center}

\textbf{Tableau de conjugaison complet} (verbe \eng{to work}) :

\begin{center}
\begin{tabular}{|l|l|l|}
\hline
\rowcolor{popblueL} \textbf{Affirmatif} & \textbf{Négatif} & \textbf{Interrogatif} \\
\hline
I will be working & I won't be working & Will I be working? \\
\hline
You will be working & You won't be working & Will you be working? \\
\hline
He/She/It will be working & He/She/It won't be working & Will he/she/it be working? \\
\hline
We will be working & We won't be working & Will we be working? \\
\hline
They will be working & They won't be working & Will they be working? \\
\hline
\end{tabular}
\end{center}

\begin{attention}[Piège n° 1 : will ne se conjugue jamais]
Les francophones, habitués à conjuguer leurs auxiliaires, commettent deux erreurs typiques :
\begin{itemize}
\item \eng{He wills be working.} — FAUX : \eng{will} ne prend jamais de \eng{-s}, même à la troisième personne. Correct : \eng{He will be working.}
\item \eng{He will is working.} / \eng{They will are working.} — FAUX : après \eng{will}, on utilise toujours la base verbale \eng{be}. Correct : \eng{He will be working.}
\end{itemize}
Retenez : \eng{will} est un modal, et les modaux anglais sont \emph{invariables} et toujours suivis de la base verbale.
\end{attention}

\begin{popfigure}
\begin{tikzpicture}[scale=1]
% Ligne du temps
\draw[-{Stealth[length=3mm]}, thick, popdark] (-0.5,0) -- (11,0) node[below left, font=\small] {temps};
% Point "now"
\fill[popgreen] (1.5,0) circle (0.12);
\node[below, font=\small\bfseries, popgreen] at (1.5,-0.15) {now};
% Point "this time tomorrow"
\fill[poporange] (7.5,0) circle (0.12);
\node[below, font=\small\bfseries, poporange, align=center] at (7.5,-0.15) {this time\\tomorrow};
% Vague représentant l'action en cours
\draw[very thick, popblue, decorate, decoration={snake, amplitude=0.18cm, segment length=0.5cm}] (5.2,0.55) -- (9.8,0.55);
\node[above, font=\small, popblue, align=center] at (7.5,0.95) {action en cours :\\ \eng{I will be sitting in the exam room}};
% Flèche vers le point
\draw[-{Stealth}, poporange, dashed] (7.5,0.45) -- (7.5,0.15);
% Futur
\node[font=\small\itshape, popmuted] at (9.8,-0.6) {futur};
\end{tikzpicture}
\legende{Le future continuous : au moment précis « this time tomorrow », l'action (la vague bleue) est en train de se dérouler.}
\end{popfigure}

\courssub{Les trois emplois du future continuous}

\textbf{Emploi 1 : une action en cours à un moment précis du futur}

C'est l'emploi fondamental. On imagine un instant précis dans le futur et l'on décrit ce qui sera \emph{en train de se passer} à cet instant-là. Les marqueurs typiques sont : \eng{this time tomorrow} (« demain à la même heure »), \eng{at 8 p.m.}, \eng{this time next week}, \eng{at this moment next month}.

\begin{exemplebox}[Exemples traduits]
\begin{itemize}
\item \eng{At 8 p.m., we will be having dinner.} « À 20 heures, nous serons en train de dîner. »
\item \eng{This time next week, we will be writing the Bac exam.} « À la même heure la semaine prochaine, nous serons en train de composer le Bac. »
\item \eng{Don't call me at noon; I will be attending a meeting.} « Ne m'appelle pas à midi ; je serai en réunion. »
\item \eng{This time next month, they will be flying to London.} « Le mois prochain à la même heure, ils seront dans l'avion pour Londres. »
\end{itemize}
\end{exemplebox}

\textbf{Emploi 2 : une action future prévue, faisant partie d'un programme normal}

Le future continuous présente l'action future comme quelque chose de \emph{naturel, déjà prévu}, qui arrivera « de toute façon », dans le cours normal des choses.

\begin{exemplebox}[Exemples traduits]
\begin{itemize}
\item \eng{I will be seeing the principal tomorrow anyway, so I can give him your message.} « De toute façon, je verrai le proviseur demain, donc je peux lui remettre ton message. »
\item \eng{She will be studying at university next year.} « Elle sera à l'université l'année prochaine. » (c'est prévu, c'est le programme)
\item \eng{We will be travelling to Buea for the holidays.} « Nous irons à Buea pour les vacances. » (voyage prévu)
\end{itemize}
\end{exemplebox}

\textbf{Emploi 3 : une question polie sur les projets de quelqu'un}

À l'interrogatif, \eng{Will you be + V-ing ?} est une manière \emph{indirecte et polie} de s'enquérir des projets d'autrui, sans mettre la personne sous pression.

\begin{exemplebox}[Exemples traduits]
\begin{itemize}
\item \eng{Will you be coming to the meeting?} « Est-ce que tu viendras à la réunion ? » (question polie)
\item \eng{Will you be using your phone tonight?} « Est-ce que tu te serviras de ton téléphone ce soir ? »
\item \eng{Will you be staying in Douala?} « Est-ce que tu resteras à Douala ? »
\end{itemize}
\end{exemplebox}

\begin{savaistu}[La politesse britannique]
\eng{Will you be + V-ing ?} est une formule de politesse très appréciée à l'oral britannique. Demander \eng{Will you come to the party?} peut sonner comme une invitation directe à laquelle il est difficile de dire non ; en revanche, \eng{Will you be coming to the party?} laisse à l'interlocuteur une porte de sortie élégante : on s'informe simplement de son « programme ». C'est la même nuance qu'entre « Tu viens ? » et « Est-ce que tu comptais venir ? » en français. Les hôtels, les restaurants et les compagnies aériennes britanniques adorent cette tournure : \eng{Will you be needing a taxi, sir?}
\end{savaistu}

\courssub{Future continuous, will simple, be going to ou present continuous ?}

Le français utilise souvent le même futur simple là où l'anglais distingue quatre constructions. Voici le tableau comparatif essentiel pour le Bac (rappel des temps de Première + nouveau temps) :

\begin{center}
\begin{tabularx}{\linewidth}{|X|X|X|X|}
\hline
\rowcolor{popblueL} \textbf{will + BV} & \textbf{be going to + BV} & \textbf{Present Continuous} & \textbf{will be + V-ing} \\
\hline
Décision spontanée, prise au moment de parler ; promesse ; prédiction d'opinion & Intention, projet déjà décidé avant de parler ; prédiction fondée sur un signe visible & Arrangement fixe, rendez-vous planifié (souvent avec indicateur de temps) & Action en cours à un moment précis du futur ; programme normal/routine ; question polie \\
\hline
\eng{I'm thirsty. I will drink some water.} & \eng{I am going to visit Bamenda next month.} & \eng{I am seeing the principal at 10 a.m. tomorrow.} & \eng{This time tomorrow, I will be visiting Bamenda.} \\
\hline
« J'ai soif. Je vais boire de l'eau. » (décision immédiate) & « Je vais visiter Bamenda le mois prochain. » (projet décidé) & « Je vois le proviseur demain à 10 h. » (rendez-vous fixé) & « Demain à cette heure, je serai en train de visiter Bamenda. » (action en cours) \\
\hline
\end{tabularx}
\end{center}

\begin{methode}[Comment choisir le bon futur]
Posez-vous ces questions dans l'ordre :
\begin{enumerate}
\item La décision est-elle prise \emph{maintenant}, au moment où je parle (promesse, offre, décision soudaine) ? → utilisez \cle{will + base verbale}. Exemple : \eng{Someone is at the door — I will open it.}
\item S'agit-il d'un \emph{projet déjà décidé} (intention) avant ce moment, ou d'une prédiction avec un signe visible ? → utilisez \cle{be going to}. Exemple : \eng{Look at those clouds! It is going to rain.}
\item S'agit-il d'un \emph{rendez-vous fixé}, d'un arrangement précis avec une heure/date ? → utilisez le \cle{Present Continuous} (révision). Exemple : \eng{We are leaving for Douala at 6 a.m. tomorrow.}
\item L'action sera-t-elle \emph{en cours à un moment précis du futur}, ou fait-elle partie d'un programme normal/routine, ou s'agit-il d'une question polie ? → utilisez \cle{will be + V-ing}. Exemple : \eng{At 10 a.m., we will be writing the test.}
\end{enumerate}
\end{methode}

\begin{attention}[Piège n° 2 : ne pas confondre will simple et future continuous]
Comparez :
\begin{itemize}
\item \eng{I will write the essay tonight.} = « J'écrirai la rédaction ce soir. » (décision, promesse : l'accent est sur le \emph{résultat} / l'accomplissement)
\item \eng{At 9 p.m., I will be writing the essay.} = « À 21 heures, je serai en train d'écrire la rédaction. » (l'accent est sur le \emph{déroulement} à cet instant)
\end{itemize}
Autre erreur fréquente : oublier \eng{be}. \eng{I will working tomorrow.} est FAUX ; il faut \eng{I will be working tomorrow.} Enfin, n'ajoutez jamais \eng{to} après \eng{will} : \eng{will to be working} n'existe pas.
\end{attention}

\courssub{Entraînement guidé}

\begin{exoresolu}[Mettez les verbes au futur approprié]
Énoncé — Complétez avec \eng{will + BV}, \eng{be going to + BV}, \eng{Present Continuous} ou \eng{will be + V-ing}, puis justifiez.

1. \eng{This time tomorrow, we (sit) in the examination hall.}
2. \eng{— Oh no, I have forgotten my pen! — Don't worry, I (lend) you mine.}
3. \eng{My sister has bought her plane ticket: she (fly) to Nairobi on Friday.}
4. \eng{(you / use) your bicycle this afternoon? May I borrow it?}
5. \eng{At 7 a.m. tomorrow, the candidates (wait) in front of the school gate.}
6. \eng{We (have) a staff meeting at 3 p.m. tomorrow; it's on the schedule.}
\tcblower
Solution détaillée :

1. \eng{This time tomorrow, we \textbf{will be sitting} in the examination hall.} — Action en cours à un moment précis du futur (\eng{this time tomorrow}) → future continuous. Attention à la graphie : \eng{sit} → \eng{sitting} (doublement du \eng{t}).

2. \eng{Don't worry, I \textbf{will lend} you mine.} — Décision spontanée, prise à l'instant où l'on parle → \eng{will} simple.

3. \eng{She \textbf{is going to fly} to Nairobi on Friday.} — Projet déjà décidé et organisé (le billet est acheté) → \eng{be going to}. (Le \eng{Present Continuous} \eng{is flying} serait aussi possible pour un horaire fixe, mais \eng{be going to} insiste sur l'intention réalisée).

4. \eng{\textbf{Will you be using} your bicycle this afternoon?} — Question polie sur les projets de quelqu'un → future continuous à l'interrogatif : \eng{Will + sujet + be + V-ing}.

5. \eng{At 7 a.m. tomorrow, the candidates \textbf{will be waiting} in front of the school gate.} — Action en cours à un moment précis du futur (\eng{at 7 a.m. tomorrow}) → future continuous.

6. \eng{We \textbf{are having} a staff meeting at 3 p.m. tomorrow; it's on the schedule.} — Arrangement fixe, horaire précis → Present Continuous (révision).
\end{exoresolu}

\begin{exercice}[À vous de jouer]
1. Traduisez en anglais : « Demain à la même heure, nous serons en train de réviser la grammaire. »
2. Traduisez en anglais, avec une question polie : « Est-ce que tu passeras au marché ce soir ? »
3. Corrigez la phrase fautive : \eng{He wills be work at 9 o'clock.}
4. Choisissez la bonne forme et justifiez : \eng{Watch out! You (drop / are going to drop / will be dropping) your phone!}
5. Complétez : \eng{Next Monday at this time, I (fly) \_\_\_\_\_\_\_\_\_\_ to London.} (Contexte : action en cours à un moment précis).
\end{exercice}

\begin{corrige}[Exercice « À vous de jouer »]
1. \eng{This time tomorrow, we will be revising grammar.} (action en cours à un moment précis du futur)
2. \eng{Will you be going to the market this evening?} (question polie : \eng{Will you be + V-ing})
3. \eng{He will be working at 9 o'clock.} — \eng{will} est invariable (pas de \eng{-s}) et il faut la base verbale \eng{be} suivie du verbe en \eng{-ing}.
4. \eng{Watch out! You \textbf{are going to drop} your phone!} — Prédiction fondée sur un signe visible (le téléphone glisse) → \eng{be going to}.
5. \eng{Next Monday at this time, I \textbf{will be flying} to London.} (Future continuous : \eng{at this time} + action en cours).
\end{corrige}

\begin{aretenir}[L'essentiel du future continuous]
\begin{itemize}
\item \textbf{Forme} : sujet + \eng{will be + V-ing} ; négation \eng{won't be + V-ing} ; question \eng{Will + sujet + be + V-ing ?}
\item \eng{will} est invariable ; \eng{be} reste à la base verbale ; le verbe principal prend \eng{-ing} (attention aux règles d'orthographe : \eng{write $\to$ writing}, \eng{sit $\to$ sitting}).
\item \textbf{Sens} : une action \emph{en train de se dérouler} à un moment précis du futur (\eng{this time tomorrow}, \eng{at 8 p.m.}).
\item \textbf{Emplois} : action en cours à un instant futur ; programme normal/routine déjà prévu ; question polie sur les projets (\eng{Will you be coming?}).
\item \textbf{Choix du futur} : décision spontanée → \eng{will} ; intention/projet → \eng{be going to} ; rendez-vous fixe → \eng{Present Continuous} ; action en cours à un moment futur → \eng{will be + V-ing}.
\end{itemize}
\end{aretenir}

\coursec{Comparaison français/anglais et pièges des francophones}

Nous savons maintenant former et employer le \cle{present perfect continuous} (\eng{have/has been + V-ing}) et le \cle{future continuous} (\eng{will be + V-ing}). Mais connaître les règles ne suffit pas : pour réussir au Bac, il faut comprendre \textbf{comment le français pense différemment} et repérer les pièges précis où les élèves francophones tombent presque tous. C'est l'objet de cette section : d'abord observer, puis comprendre, puis s'entraîner à éviter les cinq erreurs classiques.

\courssub{Le problème de fond : le français n'a pas ces temps}

\begin{textebox}[A conversation in Yaoundé]
— How long have you been learning English, Amina?

— I have been learning English since Form One. That's almost six years now.

— And have you been practising outside school?

— Yes, I have been watching English films every weekend, and my uncle in Buea has been sending me novels. What about you, what will you be doing during the holidays?

— This time next month, I will be travelling to Douala with my family. We will be staying with my grandmother for two weeks.
\end{textebox}

\textbf{Glossaire :} \emph{to practise} (v., s'entraîner — orthographe britannique ; le nom est \emph{practice}) ; \emph{outside school} (loc., en dehors de l'école) ; \emph{to send} (v., envoyer — irrégulier : \emph{send, sent, sent}) ; \emph{to stay} (v., séjourner).

\textbf{Observation.} Traduisons mot à mot les phrases clés :

\begin{center}
\begin{tabularx}{\linewidth}{|X|X|}
\hline
\rowcolor{popblueL} English & Traduction française naturelle \\
\hline
\eng{I have been learning English since Form One.} & J'apprends l'anglais depuis la classe de sixième. \\
\hline
\eng{It has been raining since this morning.} & Il pleut depuis ce matin. \\
\hline
\eng{She has been working here for ten years.} & Elle travaille ici depuis dix ans. \\
\hline
\eng{This time tomorrow, I will be travelling.} & À la même heure demain, je serai en train de voyager. \\
\hline
\end{tabularx}
\end{center}

Le constat est frappant : là où l'anglais emploie un temps \emph{composé et continu}, le français emploie un simple \textbf{présent} accompagné de « depuis ». Et là où l'anglais emploie naturellement le futur continu, le français se contente souvent du futur simple, la tournure « être en train de » étant rare au futur.

\begin{attention}[« Depuis » + présent français = present perfect en anglais]
C'est le piège numéro un des francophones. Quand le français dit « Il pleut \textbf{depuis} ce matin » avec un présent, l'anglais dit obligatoirement \eng{It has been raining since this morning} — jamais \eng{It rains since this morning}. Règle pratique : dès que vous voyez « depuis » ou « cela fait... que » en français, traduisez par le \cle{present perfect} (simple ou continuous), jamais par le présent simple. Attention à la distinction : \eng{since} + point de départ (\eng{since 2018}, \eng{since Monday}) ; \eng{for} + durée (\eng{for six years}, \eng{for two hours}).
\end{attention}

\begin{propriete}[Correspondances fondamentales français / anglais]
\begin{itemize}
\item « Depuis » + présent (action qui dure encore) $\rightarrow$ \textbf{present perfect continuous} : \eng{I have been living in Yaoundé since 2019.} (J'habite à Yaoundé depuis 2019.)
\item « Depuis » + présent avec un verbe d'état (know, like, have) $\rightarrow$ \textbf{present perfect simple} : \eng{I have known her since childhood.} (Je la connais depuis l'enfance.)
\item « Être en train de » au futur $\rightarrow$ \textbf{future continuous} : \eng{I will be working at 8 a.m.} (Je serai en train de travailler à 8 h.)
\end{itemize}
\end{propriete}

\courssub{Les cinq erreurs classiques des francophones}

\textbf{Erreur 1 : oublier \cle{been}.} Beaucoup d'élèves écrivent \eng{I have studying} par calque du français « j'ai étudié ». Or le present perfect continuous exige \textbf{deux} auxiliaires : \eng{have/has} + \eng{been}, puis le verbe en -ing. Correct : \eng{I have been studying for three hours.} (J'étudie depuis trois heures.) N'oubliez pas non plus l'accord : \eng{has} avec \eng{he/she/it}, \eng{have} partout ailleurs.

\textbf{Erreur 2 : utiliser le present simple avec \eng{since/for}.} La phrase \eng{I learn English since 2018} est la faute la plus fréquente au Cameroun. En anglais, \eng{since} et \eng{for} (avec une durée qui continue jusqu'à maintenant) appellent le present perfect : \eng{I have been learning English since 2018.} Comparez : \eng{He has been playing football for two hours.} (Il joue au football depuis deux heures) et non \eng{He plays football since two hours}.

\textbf{Erreur 3 : mettre un verbe d'état à la forme continue.} Les verbes d'état (\eng{know, want, like, love, hate, believe, understand, own, need, seem}) ne prennent pas la forme continue. On ne dit jamais \eng{I am knowing} ni \eng{I have been wanting}. On dit : \eng{I have known him for years.} (Je le connais depuis des années) ; \eng{She has wanted a bicycle for months.} (Elle veut un vélo depuis des mois.)

\textbf{Erreur 4 : double marqueur du futur ou oubli de \eng{be}.} Deux fautes symétriques : \eng{I will to be going} (deux marqueurs mélangés) et \eng{I will studying} (oubli de \eng{be}). La seule structure correcte est \eng{will be + V-ing} : \eng{I will be studying at this time tomorrow.}

\begin{propriete}[Un seul marqueur de futur par proposition]
En anglais, une proposition ne contient qu'\textbf{un seul} marqueur de futur : \eng{will} \textbf{OU} \eng{be going to}, jamais les deux ensemble. On écrit \eng{I will be travelling} ou \eng{I am going to travel}, mais jamais \eng{I will to be going}. De même, après \eng{will}, le verbe reste à la base verbale : \eng{will be}, jamais \eng{will to be} ni \eng{will being}.
\end{propriete}

\textbf{Erreur 5 : l'orthographe du -ing.} Trois règles à revoir : (a) un verbe en -e muet perd son e : \eng{write} $\rightarrow$ \eng{writing} (jamais \eng{writting}), \eng{make} $\rightarrow$ \eng{making} (jamais \eng{makeing}) ; (b) une consonne finale après une voyelle brève accentuée double : \eng{run} $\rightarrow$ \eng{running}, \eng{sit} $\rightarrow$ \eng{sitting}, \eng{begin} $\rightarrow$ \eng{beginning} ; (c) -ie devient -y : \eng{lie} $\rightarrow$ \eng{lying}, \eng{die} $\rightarrow$ \eng{dying}.

\begin{exemplebox}[Exemples corrects traduits]
\begin{itemize}
\item \eng{She has been working here for ten years.} — Elle travaille ici depuis dix ans.
\item \eng{I will be travelling this time tomorrow.} — Je serai en train de voyager à la même heure demain.
\item \eng{They have been building the bridge since January.} — Ils construisent le pont depuis janvier.
\item \eng{We will be sitting our exam at 8 o'clock.} — Nous serons en train de composer à 8 heures.
\item \eng{He has been running, so he is tired.} — Il vient de courir, donc il est fatigué.
\end{itemize}
\end{exemplebox}

\begin{popfigure}
\begin{center}
\begin{tabularx}{\linewidth}{|X|X|}
\hline
\rowcolor{popblueL} Phrase fautive typique & Phrase correcte + traduction \\
\hline
\emph{Faute :} I have studying English. & \eng{I have been studying English.} (J'étudie l'anglais depuis un moment.) \\
\hline
\emph{Faute :} I learn English since 2018. & \eng{I have been learning English since 2018.} (J'apprends l'anglais depuis 2018.) \\
\hline
\emph{Faute :} It rains since this morning. & \eng{It has been raining since this morning.} (Il pleut depuis ce matin.) \\
\hline
\emph{Faute :} I am knowing the answer. & \eng{I know the answer.} (Je connais la réponse.) \\
\hline
\emph{Faute :} I have been wanting this. & \eng{I have wanted this for a long time.} (Je le veux depuis longtemps.) \\
\hline
\emph{Faute :} I will to be going home. & \eng{I will be going home.} (Je serai en train de rentrer.) \\
\hline
\emph{Faute :} I will studying tonight. & \eng{I will be studying tonight.} (Je serai en train d'étudier ce soir.) \\
\hline
\emph{Faute :} She has been writting. & \eng{She has been writing.} (Elle est en train d'écrire depuis un moment.) \\
\hline
\emph{Faute :} They are makeing noise. & \eng{They are making noise.} (Ils font du bruit.) \\
\hline
\end{tabularx}
\end{center}
\legende{Les fautes les plus fréquentes des candidats francophones et leurs corrections.}
\end{popfigure}

\begin{popfigure}
\begin{center}
\begin{tikzpicture}[node distance=0.6cm,
  q/.style={draw=popblue, fill=popblueL, rounded corners, align=center, font=\small, text width=4.2cm},
  a/.style={draw=popgreen, fill=popgreenL, rounded corners, align=center, font=\small, text width=4.6cm},
  arr/.style={-{Stealth}, thick, popdark}]
\node[q] (q1) {1. Y a-t-il \textbf{been} après have/has ?};
\node[q, below=of q1] (q2) {2. Le verbe est-il bien en \textbf{-ing} ?};
\node[q, below=of q2] (q3) {3. Est-ce un \textbf{verbe d'état} ? (know, want, like...)};
\node[q, below=of q3] (q4) {4. Y a-t-il \textbf{be} après will ?};
\node[a, right=1.2cm of q2] (ok) {Si oui partout : la phrase est correcte !};
\draw[arr] (q1) -- (q2);
\draw[arr] (q2) -- (q3);
\draw[arr] (q3) -- (q4);
\draw[arr, popgreen] (q3.east) -- node[above, font=\scriptsize]{non} (ok.south west);
\end{tikzpicture}
\end{center}
\legende{Les quatre questions de vérification avant de rendre sa copie.}
\end{popfigure}

\begin{methode}[Vérifier sa copie en quatre questions]
Avant de rendre toute production écrite contenant ces temps, posez-vous dans l'ordre :
\begin{enumerate}
\item Y a-t-il \cle{been} après \eng{have/has} ? (\eng{I have been waiting}, pas \eng{I have waiting}.)
\item Le verbe principal est-il bien en \cle{-ing}, avec la bonne orthographe ? (\eng{writing}, \eng{running}, \eng{making}.)
\item Le verbe est-il un \cle{verbe d'état} ? Si oui, revenez à la forme simple (\eng{I have known}, pas \eng{I have been knowing}).
\item Pour le futur, y a-t-il \cle{be} après \eng{will}, et un seul marqueur de futur ? (\eng{I will be going}, pas \eng{I will going} ni \eng{I will to be going}.)
\end{enumerate}
\end{methode}

\begin{exoresolu}[Corrige les erreurs]
Énoncé. Chaque phrase contient une erreur. Corrige-la et traduis la phrase correcte.
\begin{enumerate}
\item \eng{My father has working in Douala for twenty years.}
\item \eng{We learn French since primary school.}
\item \eng{I am understanding the lesson now.}
\item \eng{She will to be travelling to Bamenda tomorrow.}
\item \eng{They have been siting here since 7 o'clock.}
\end{enumerate}
\tcblower
Solution détaillée.
\begin{enumerate}
\item Oubli de \eng{been} : \eng{My father has \textbf{been} working in Douala for twenty years.} — Mon père travaille à Douala depuis vingt ans.
\item Present simple interdit avec \eng{since} : \eng{We \textbf{have been learning} French since primary school.} — Nous apprenons le français depuis l'école primaire.
\item Verbe d'état à la forme continue : \eng{I \textbf{understand} the lesson now.} — Je comprends la leçon maintenant.
\item Double marqueur du futur : \eng{She will \textbf{be travelling} to Bamenda tomorrow.} — Elle sera en train de voyager vers Bamenda demain.
\item Orthographe du -ing (consonne double) : \eng{They have been \textbf{sitting} here since 7 o'clock.} — Ils sont assis ici depuis 7 heures.
\end{enumerate}
\end{exoresolu}

\begin{savaistu}[Ce que les correcteurs du Bac sanctionnent en priorité]
Au baccalauréat, les correcteurs repèrent presque systématiquement deux fautes : l'oubli de \cle{been} dans le present perfect continuous et le mauvais accord \eng{have/has} (\eng{he have} au lieu de \eng{he has}). Relire uniquement ces deux points dans votre composition suffit souvent à gagner un ou deux points précieux. Les meilleurs candidats gardent aussi une petite liste mentale des verbes d'état pour ne jamais les mettre en -ing.
\end{savaistu}

\begin{exercice}[À toi de jouer]
Traduis en anglais : 1) J'attends le bendskin depuis vingt minutes. 2) Elle connaît mon frère depuis l'enfance. 3) À la même heure demain, nous serons en train de réviser. 4) Il pleut depuis hier soir. 5) Ils construisent cette école depuis mars.
\end{exercice}

\begin{corrige}[Exercice « À toi de jouer »]
1) \eng{I have been waiting for the bendskin for twenty minutes.} 2) \eng{She has known my brother since childhood.} (verbe d'état : forme simple). 3) \eng{This time tomorrow, we will be revising.} 4) \eng{It has been raining since last night.} 5) \eng{They have been building this school since March.}
\end{corrige}

\coursec{Entraînement guidé et production}

Nous avons vu dans la section précédente comment le français et l'anglais expriment différemment la durée et le futur, et quels pièges attendent le francophone (\eng{I am living here since 2020} est impossible !). Il est temps maintenant de \cle{s'entraîner} : c'est en conjuguant, en transformant et en corrigeant que ces temps deviendront des réflexes. Cette section vous propose une méthode d'examen, cinq exercices progressifs entièrement corrigés, puis une production écrite finale.

\courssub{La méthode du marqueur temporel}

\begin{methode}[Réussir les exercices de conjugaison au Bac]
Au Bac, les exercices de grammaire demandent souvent de mettre le verbe entre parenthèses au temps correct. Procédez toujours ainsi :
\begin{enumerate}
\item \textbf{Soulignez d'abord le marqueur temporel} de la phrase : \eng{for}, \eng{since}, \eng{all morning}, \eng{this time tomorrow}, \eng{at 8 p.m. on Friday}, \eng{how long}. Il indique presque toujours le temps attendu.
\item \textbf{Posez-vous la question du sens} : l'action est-elle terminée avec un résultat visible (\textrightarrow{} present perfect simple) ? Est-elle encore en cours ou présentée comme une durée (\textrightarrow{} present perfect continuous) ? Se déroulera-t-elle à un moment précis du futur (\textrightarrow{} future continuous) ?
\item \textbf{Construisez la forme} : sujet + \eng{have/has been} + verbe-\eng{ing} (present perfect continuous) ; sujet + \eng{will be} + verbe-\eng{ing} (future continuous).
\item \textbf{Vérifiez l'accord} (\eng{has} avec he/she/it, \eng{have} ailleurs) et l'orthographe du participe présent (\eng{study} \textrightarrow{} \eng{studying}, \eng{write} \textrightarrow{} \eng{writing}, \eng{run} \textrightarrow{} \eng{running}).
\item \textbf{Relisez la phrase entière} pour contrôler la cohérence : un repère passé daté (\eng{yesterday}, \eng{last week}, \eng{in 2019}) interdit le present perfect et impose le preterit simple.
\end{enumerate}
\end{methode}

\begin{popfigure}
\begin{tikzpicture}[
node distance=0.9cm and 1.4cm,
startstop/.style={rectangle, rounded corners, draw=popblue, fill=popblueL, text=popdark, align=center, minimum height=0.9cm, inner sep=4pt},
decision/.style={diamond, draw=poporange, fill=poporangeL, text=popdark, align=center, aspect=2.2, inner sep=1pt},
result/.style={rectangle, draw=popgreen, fill=popgreenL, text=popdark, align=center, minimum height=0.9cm, inner sep=4pt},
arrow/.style={-{Stealth[length=2.5mm]}, thick, popdark}]
\node[startstop, align=center] (start){Phrase à conjuguer :\\repérer le marqueur temporel};
\node[decision, below=of start, align=center] (q1){Action terminée\\avec résultat ?};
\node[decision, below left=1.1cm and 0.6cm of q1, align=center] (q2){Durée ou\\processus ?};
\node[decision, below right=1.1cm and 0.6cm of q1, align=center] (q3){Moment précis\\du futur ?};
\node[result, below=of q2, align=center] (r1){\textbf{Present perfect}\\\textbf{continuous}\\have/has been + V-ing};
\node[result, below=of q3, align=center] (r2){\textbf{Future continuous}\\will be + V-ing};
\node[result, right=1.2cm of r2, align=center] (r3){\textbf{Present perfect simple}\\have/has + participe passé};
\node[result, left=1.2cm of r1, align=center] (r4){\textbf{Future simple}\\will + base verbale};
\draw[arrow] (start) -- (q1);
\draw[arrow] (q1) -| node[pos=0.25, above left, font=\small]{non} (q2);
\draw[arrow] (q1) -| node[pos=0.25, above right, font=\small]{oui} (r3.north);
\draw[arrow] (q2) -- node[left, font=\small]{oui} (r1);
\draw[arrow] (q3) -- node[right, font=\small]{oui} (r2);
\draw[arrow] (q2.east) -- node[above, font=\small]{non} (q3.west);
\draw[arrow] (q3.south) |- node[pos=0.75, above, font=\small]{non} (r4.east);
\end{tikzpicture}
\legende{Organigramme de décision : choisir entre present perfect simple, present perfect continuous, future continuous et future simple.}
\end{popfigure}

\courssub{Exercices résolus : les deux exemples de référence}

\begin{exoresolu}[Exemple traité 1 — durée, action non terminée]
Énoncé : mettez le verbe entre parenthèses au temps correct. \eng{She (study) medicine since 2020.}
\tcblower
Solution détaillée :
\begin{enumerate}
\item Marqueur temporel : \eng{since 2020} \textrightarrow{} action commencée dans le passé et qui continue.
\item Sens : elle étudie encore la médecine aujourd'hui ; on insiste sur la durée, pas sur un résultat.
\item Forme : sujet \eng{she} \textrightarrow{} \eng{has been} + \eng{studying}.
\end{enumerate}
Réponse : \eng{She has been studying medicine since 2020.} « Elle étudie la médecine depuis 2020. »
\end{exoresolu}

\begin{exoresolu}[Exemple traité 2 — action en cours à un moment du futur]
Énoncé : \eng{At 10 a.m. tomorrow, we (write) the English paper.}
\tcblower
Solution détaillée :
\begin{enumerate}
\item Marqueur temporel : \eng{at 10 a.m. tomorrow} \textrightarrow{} moment précis du futur.
\item Sens : à 10 heures précises, l'action sera en train de se dérouler.
\item Forme : \eng{will be} + \eng{writing} (attention : \eng{write} perd son \eng{-e} final).
\end{enumerate}
Réponse : \eng{At 10 a.m. tomorrow, we will be writing the English paper.} « Demain à 10 heures, nous serons en train de composer l'épreuve d'anglais. »
\end{exoresolu}

\courssub{Exercice 1 : present perfect simple ou continuous ?}

\begin{exercice}[Complétez avec le present perfect simple ou continuous]
Mettez le verbe entre parenthèses au present perfect simple ou continuous. Justifiez votre choix (résultat ou durée ?).
\begin{enumerate}
\item I (read) this novel for three weeks and I still have fifty pages left.
\item My father (repair) the car at last! We can travel to Douala now.
\item How long (you / learn) English?
\item The students (write) three essays this term.
\item It (rain) since six o'clock this morning; the road to Bamenda is flooded.
\item Amina (not / finish) her homework yet.
\item We (wait) for the bus for forty minutes!
\item Someone (steal) my bicycle!
\end{enumerate}
\end{exercice}

\begin{corrige}[Exercice 1]
\begin{enumerate}
\item \eng{I have been reading this novel for three weeks...} — durée (\eng{for three weeks}), action non terminée.
\item \eng{My father has repaired the car at last!} — résultat visible : la voiture marche.
\item \eng{How long have you been learning English?} — question de durée avec \eng{how long}.
\item \eng{The students have written three essays this term.} — bilan chiffré, résultat.
\item \eng{It has been raining since six o'clock this morning...} — durée, action en cours.
\item \eng{Amina has not finished her homework yet.} — résultat attendu, avec \eng{yet}.
\item \eng{We have been waiting for the bus for forty minutes!} — durée, irritation.
\item \eng{Someone has stolen my bicycle!} — résultat : le vélo a disparu (on ne peut pas « voler en continu » un objet précis).
\end{enumerate}
\end{corrige}

\courssub{Exercice 2 : le future continuous à partir d'un emploi du temps}

\begin{exercice}[Demain, à cette heure-là...]
Voici l'emploi du temps de Chantal, élève en Tle A à Buea, pour demain. Rédigez une phrase au future continuous pour chaque horaire, selon le modèle.
Modèle : 7.00 — have breakfast \textrightarrow{} \eng{At 7 o'clock tomorrow, Chantal will be having breakfast.}
\begin{enumerate}
\item 8.00 — attend the Maths lesson
\item 10.30 — revise in the school library
\item 12.00 — have lunch with her classmates
\item 15.00 — play football on the school field
\item 18.00 — travel back home by bendskin
\item 20.00 — do her English homework
\end{enumerate}
\end{exercice}

\begin{corrige}[Exercice 2]
\begin{enumerate}
\item \eng{At 8 o'clock, she will be attending the Maths lesson.}
\item \eng{At 10.30, she will be revising in the school library.}
\item \eng{At 12 o'clock, she will be having lunch with her classmates.}
\item \eng{At 3 p.m., she will be playing football on the school field.}
\item \eng{At 6 p.m., she will be travelling back home by bendskin.} (orthographe britannique : \eng{travel} \textrightarrow{} \eng{travelling}, avec doublement du \eng{-l})
\item \eng{At 8 p.m., she will be doing her English homework.}
\end{enumerate}
Remarque : toutes ces phrases présentent l'action comme \cle{en cours} à un moment précis du futur — c'est l'emploi fondamental du future continuous.
\end{corrige}

\courssub{Exercice 3 : transformation et changement de sens}

\begin{exercice}[Du simple au continuous]
Quand c'est possible, transformez la phrase en present perfect continuous et expliquez le changement de sens. Quand la transformation est impossible, dites pourquoi.
\begin{enumerate}
\item I have read three chapters of the book.
\item She has known her best friend since childhood.
\item They have played football all afternoon.
\item He has broken his leg.
\item We have lived in Yaoundé for ten years.
\end{enumerate}
\end{exercice}

\begin{corrige}[Exercice 3]
\begin{enumerate}
\item \eng{I have been reading the book.} — le simple insiste sur le résultat (trois chapitres terminés) ; le continuous insiste sur l'activité en cours, sans préciser la quantité lue.
\item Transformation impossible : \eng{know} est un verbe d'état ; il n'admet pas la forme continue. On garde \eng{She has known her best friend since childhood.}
\item \eng{They have been playing football all afternoon.} — le simple constate le fait ; le continuous souligne la durée (\eng{all afternoon}) et suggère qu'ils jouent peut-être encore.
\item Transformation impossible : \eng{break} exprime ici un événement instantané avec un résultat (la jambe est cassée). On garde \eng{He has broken his leg.}
\item \eng{We have been living in Yaoundé for ten years.} — les deux formes sont correctes ; le continuous met légèrement plus l'accent sur la durée vécue. (Avec \eng{live} et \eng{work}, les deux formes sont pratiquement équivalentes.)
\end{enumerate}
\end{corrige}

\begin{attention}[How long : le present perfect est obligatoire]
Dans les questions avec \eng{How long...?} portant sur une situation qui continue, le \cle{present perfect} (simple ou continuous) est obligatoire en anglais, alors que le français utilise le présent :
\begin{itemize}
\item \eng{How long have you been living here?} « Depuis combien de temps habites-tu ici ? » — et non \eng{How long do you live here?}
\item \eng{How long has she been studying medicine?} « Depuis combien de temps étudie-t-elle la médecine ? »
\end{itemize}
De même, on ne dit jamais \eng{I am living here since 2020} mais \eng{I have been living here since 2020}. Le présent simple ou continu + \eng{since/for} est l'erreur de francophone la plus fréquente au Bac.
\end{attention}

\courssub{Exercice 4 : correction de phrases d'élèves francophones}

\begin{exercice}[Trouvez et corrigez l'erreur]
Chaque phrase ci-dessous a été écrite par un élève francophone et contient une erreur. Corrigez-la et expliquez.
\begin{enumerate}
\item I am waiting for you since two hours.
\item Tomorrow at 9 a.m., I will write my exam.
\item How long do you know your English teacher?
\item She has been knowing him for years.
\item We have been finishing our project last week.
\end{enumerate}
\end{exercice}

\begin{corrige}[Exercice 4]
\begin{enumerate}
\item \eng{I have been waiting for you for two hours.} — durée commencée dans le passé et qui continue : present perfect continuous ; \eng{since} + point de départ, \eng{for} + durée.
\item \eng{Tomorrow at 9 a.m., I will be writing my exam.} — action en cours à un moment précis du futur : future continuous.
\item \eng{How long have you known your English teacher?} — \eng{how long} exige le present perfect ; \eng{know} est un verbe d'état, donc forme simple.
\item \eng{She has known him for years.} — \eng{know} ne se met jamais à la forme continue.
\item \eng{We finished our project last week.} — \eng{last week} est un repère passé terminé : preterit simple ; le present perfect est incompatible avec un moment passé daté.
\end{enumerate}
\end{corrige}

\courssub{Exercice 5 : traduction (thème)}

\begin{exercice}[Traduisez en anglais]
\begin{enumerate}
\item J'apprends l'anglais depuis six ans.
\item Elle prépare le concours depuis le mois de janvier.
\item Demain à cette heure-ci, nous serons en train de voyager vers Douala.
\item Depuis combien de temps travaillez-vous dans cette école ?
\item À 8 heures ce soir, ils regarderont le match de football.
\end{enumerate}
\end{exercice}

\begin{corrige}[Exercice 5]
\begin{enumerate}
\item \eng{I have been learning English for six years.}
\item \eng{She has been preparing for the competitive examination since January.}
\item \eng{This time tomorrow, we will be travelling to Douala.}
\item \eng{How long have you been working in this school?}
\item \eng{At 8 o'clock tonight, they will be watching the football match.}
\end{enumerate}
Observez : le présent français des phrases 1, 2 et 4 devient un present perfect continuous en anglais ; « être en train de » au futur (phrase 3) se rend par le future continuous.
\end{corrige}

\begin{center}
\begin{tabularx}{\linewidth}{|X|X|X|}
\hline
\rowcolor{popblueL} Français & English & Remarque \\
\hline
J'attends depuis une heure. & I have been waiting for an hour. & présent français \textrightarrow{} present perfect continuous \\
\hline
Depuis combien de temps habites-tu ici ? & How long have you been living here? & \eng{how long} + present perfect obligatoire \\
\hline
Demain à 10 h, je serai en train de composer. & At 10 a.m. tomorrow, I will be writing my exam. & action en cours \textrightarrow{} future continuous \\
\hline
Il pleut depuis ce matin. & It has been raining since this morning. & \eng{since} + point de départ \\
\hline
Elle connaît Paul depuis des années. & She has known Paul for years. & verbe d'état : forme simple \\
\hline
\end{tabularx}
\end{center}

\courssub{Production finale : My life now and my future}

\begin{experience}[Rédaction guidée — My life now and my future]
Rédigez un court paragraphe (80 à 120 mots) intitulé \eng{My life now and my future}. Contraintes obligatoires :
\begin{itemize}
\item au moins \textbf{deux present perfect continuous} (activités que vous faites depuis un certain temps) ;
\item au moins \textbf{deux future continuous} (ce que vous serez en train de faire à des moments précis) ;
\item au moins un marqueur \eng{for}, un \eng{since}, un \eng{this time + futur} et un \eng{at + heure}.
\end{itemize}
Plan suggéré : (1) votre vie d'élève actuelle et sa durée ; (2) vos projets et ce que vous ferez l'année prochaine à un moment précis.
\end{experience}

\begin{exemplebox}[Modèle de production]
\eng{My life now and my future — I have been studying at Government High School Buea for six years, and I have been preparing seriously for the Baccalaureate since September. English has become very important to me because I have been reading English novels every evening. This time next year, I will be sitting in a university lecture hall in Yaoundé, and I will be studying medicine, which has always been my dream. At 8 p.m. tonight, I will be revising my grammar lessons with my classmates. I am sure all this hard work will open many doors for me.} (102 mots)

« Ma vie maintenant et mon avenir — J'étudie au lycée de Buea depuis six ans et je prépare sérieusement le baccalauréat depuis septembre. L'anglais est devenu très important pour moi parce que je lis des romans anglais chaque soir. À la même époque l'année prochaine, je serai assis dans un amphithéâtre de l'université de Yaoundé et j'étudierai la médecine, qui a toujours été mon rêve. Ce soir à 20 heures, je réviserai mes leçons de grammaire avec mes camarades. Je suis sûr que tout ce travail m'ouvrira de nombreuses portes. »
\end{exemplebox}

\begin{savaistu}[Un atout pour le Bac]
Dans les épreuves de rédaction du Baccalauréat, les correcteurs valorisent les copies qui montrent une \cle{maîtrise variée des temps}. Alterner naturellement un present perfect continuous (\eng{I have been working on this project for months}) et un future continuous (\eng{Next year, I will be studying at university}) prouve que vous dépassez le simple présent et le futur simple. Deux ou trois emplois corrects de ces temps dans une composition peuvent faire la différence entre une note moyenne et une excellente note. Entraînez-vous donc à les glisser dans toutes vos productions écrites !
\end{savaistu}

\newpage

\coursec{Activité d'intégration}

\courssub{Objectifs de l'activité}

À l'issue de cette activité d'intégration, tu seras capable de :

\begin{itemize}
\item mener un \cle{échange oral} authentique en anglais en posant des questions et en réagissant aux réponses de ton partenaire ;
\item employer correctement le \cle{present perfect continuous} (\eng{have/has been + V-ing}) pour parler d'actions commencées dans le passé et qui durent encore ;
\item employer correctement le \cle{future continuous} (\eng{will be + V-ing}) pour parler d'actions en cours à un moment précis du futur ;
\item éviter les erreurs typiques des francophones : oubli de \eng{been}, emploi du present perfect continuous avec un verbe d'état, traduction mot à mot de « depuis » ;
\item rédiger un \cle{court compte rendu écrit} à la troisième personne, en accordant correctement l'auxiliaire (\eng{has}, \eng{he will be}) ;
\item repérer et corriger des erreurs de grammaire dans une production écrite.
\end{itemize}

\courssub{Situation}

Le club d'anglais de ton lycée à Yaoundé prépare son magazine de fin d'année, intitulé \emph{The Horizon Magazine}. Le thème du prochain numéro est : \eng{“Ten Years From Now: Where Will We Be?”}. Chaque membre du club doit interviewer un camarade sur sa vie actuelle (ce qu'il fait depuis longtemps) et sur ses projets (ce qu'il sera en train de faire dans le futur), puis rédiger un court portrait qui sera publié dans le magazine.

Voici le texte d'annonce publié par le président du club :

\begin{textebox}[English Club Notice — The Horizon Magazine]
Dear members,

Our next issue will be called “Ten Years From Now”. We want to know who you are today and who you will become tomorrow.

Your task is simple: interview a classmate. Ask him or her how long he or she has been living in the neighbourhood, how long he or she has been studying English, and what he or she will be doing this time next year and in ten years.

Then write a short portrait (about ten lines) for the magazine. Begin like this: “My classmate has been living in... In ten years, he or she will be...”

Be careful with your tenses! The best portraits will be read aloud during our closing ceremony in June.

The President
\end{textebox}

\begin{definition}[Glossaire]
\begin{itemize}
\item \eng{issue} (n.) : numéro (d'un magazine).
\item \eng{task} (n.) : tâche.
\item \eng{neighbourhood} (n.) : quartier.
\item \eng{portrait} (n.) : portrait, courte présentation écrite.
\item \eng{aloud} (adv.) : à voix haute.
\item \eng{closing ceremony} (n.) : cérémonie de clôture.
\end{itemize}
\end{definition}

\courssub{Consignes}

\textbf{Tâche 1 — Préparation de l'interview (travail individuel, 5 minutes).} Lis attentivement la notice du club. Recopie les quatre questions imposées ci-dessous et vérifie que tu comprends chacune d'elles. Souligne les auxiliaires et les formes verbales.

\begin{enumerate}
\item \eng{How long have you been living in your neighbourhood?}
\item \eng{How long have you been studying English?}
\item \eng{What will you be doing this time next year?}
\item \eng{Where will you be studying or working in ten years?}
\end{enumerate}

\textbf{Tâche 2 — L'interview (par binômes, 10 minutes).} Interviewe ton partenaire avec les quatre questions. Ton partenaire répond par des phrases complètes. Prends des notes en anglais. Puis inversez les rôles.

\begin{attention}[Pendant l'interview]
Réponds toujours par une phrase complète, pas par un seul mot. Dis \eng{I have been living in Mvog-Ada for twelve years.} et non simplement \eng{Twelve years}. Attention à la prononciation de \eng{been} : « biine », et non « bine ».
\end{attention}

\textbf{Tâche 3 — Rédaction du portrait (travail individuel, 15 minutes).} À partir de tes notes, rédige un court paragraphe (8 à 12 lignes) de présentation de ton camarade, à la troisième personne. Ton portrait doit contenir :

\begin{enumerate}
\item au moins \cle{deux phrases au present perfect continuous} avec \eng{for} ou \eng{since} ;
\item au moins \cle{deux phrases au future continuous} avec un marqueur de temps futur (\eng{this time next year}, \eng{in ten years}, \eng{by 2035}) ;
\item une phrase de conclusion personnelle (\eng{I think...}, \eng{I am sure that...}).
\end{enumerate}

\textbf{Tâche 4 — Lecture à la classe (5 minutes par binôme).} Lis ton portrait à la classe. Les autres élèves écoutent et relèvent une information qu'ils trouvent intéressante.

\textbf{Tâche 5 — Correction collective (10 minutes).} L'enseignant relève au tableau les erreurs typiques entendues pendant les lectures (oubli de \eng{been}, verbes d'état au present perfect continuous, oubli de \eng{be} au future continuous). La classe corrige collectivement chaque phrase fautive.

\courssub{Ressources à mobiliser}

\begin{propriete}[Rappel : le present perfect continuous]
\textbf{Forme} : sujet + \eng{have/has been} + verbe en \eng{-ing}.\\
\textbf{Emploi} : une action commencée dans le passé et qui continue encore (souvent avec \eng{for} + durée ou \eng{since} + point de départ), ou qui vient de se terminer avec un résultat visible.
\begin{itemize}
\item \eng{She has been studying English for six years.} « Elle étudie l'anglais depuis six ans. »
\item \eng{They have been playing football since four o'clock.} « Ils jouent au football depuis quatre heures. »
\end{itemize}
\textbf{Exception} : les verbes d'état (\eng{know}, \eng{like}, \eng{want}, \eng{believe}, \eng{have} = posséder) ne prennent pas la forme continue : \eng{I have known him for years.} (et non \emph{I have been knowing}).
\end{propriete}

\begin{propriete}[Rappel : le future continuous]
\textbf{Forme} : sujet + \eng{will be} + verbe en \eng{-ing}.\\
\textbf{Emplois} : (1) une action en cours à un moment précis du futur ; (2) une action prévue, qui fera partie du déroulement normal des choses.
\begin{itemize}
\item \eng{This time next year, I will be writing my A-Level exams.} « À la même époque l'année prochaine, je serai en train de passer mes examens du A-Level. »
\item \eng{In ten years, she will be working as a doctor in Douala.} « Dans dix ans, elle travaillera comme médecin à Douala. »
\end{itemize}
\end{propriete}

\begin{attention}[Pièges des francophones]
\begin{itemize}
\item Ne traduis jamais « depuis » par \eng{since} + durée : on dit \eng{for six years} (durée) mais \eng{since 2019} (point de départ).
\item N'oublie pas \eng{been} : \eng{I have been living...} et non \emph{I have living}.
\item N'oublie pas \eng{be} au future continuous : \eng{I will be studying} et non \emph{I will studying}.
\item Le français « j'habite ici depuis dix ans » utilise le présent ; l'anglais exige le present perfect continuous : \eng{I have been living here for ten years.}
\end{itemize}
\end{attention}

\begin{center}
\begin{tabularx}{\linewidth}{|X|X|X|}
\hline
\rowcolor{popblueL} Français & English & Remarque \\
\hline
J'étudie l'anglais depuis six ans. & I have been studying English for six years. & « depuis » + durée = \eng{for} \\
\hline
Il habite à Buea depuis 2018. & He has been living in Buea since 2018. & « depuis » + date = \eng{since} \\
\hline
Je le connais depuis longtemps. & I have known him for a long time. & verbe d'état : pas de forme continue \\
\hline
L'année prochaine à cette époque, je passerai mon bac. & This time next year, I will be sitting my A-Levels. & \eng{will be} + V-ing \\
\hline
\end{tabularx}
\end{center}

\textbf{Lexique utile} : \eng{neighbourhood} (quartier), \eng{secondary school} (collège/lycée), \eng{university} (université), \eng{to run a business} (gérer une entreprise), \eng{nurse} (infirmier/infirmière), \eng{engineer} (ingénieur), \eng{to dream of} (rêver de), \eng{abroad} (à l'étranger), \eng{market} (marché).

\courssub{Proposition de production et corrigé commenté}

\begin{exoresolu}[Portrait modèle pour The Horizon Magazine]
Rédige le portrait de ton camarade interviewé (8 à 12 lignes) en respectant les consignes de la tâche 3.
\tcblower
\textbf{Production modèle :}

\eng{My classmate is called Sandra. She has been living in the Nkolbisson neighbourhood for fifteen years, since she was a small child. She has been studying English since primary school, so for almost ten years now, and she speaks it very well. She has also been playing basketball for the school team for two years. This time next year, she will be preparing for her university entrance exams, because she wants to study medicine. In ten years, she will be working as a doctor, probably in a big hospital in Douala or Yaoundé. She says she will be living near her family, because family is very important to her. I am sure Sandra will achieve her dreams, because she is hardworking and serious.}

\textbf{Commentaire des choix de langue :}
\begin{itemize}
\item \eng{has been living... for fifteen years} : present perfect continuous avec \eng{has} (troisième personne) et \eng{for} + durée ; \eng{live} n'est pas un verbe d'état, la forme continue est possible.
\item \eng{since she was a small child} : \eng{since} + point de départ, suivi d'une proposition au prétérit simple.
\item \eng{has been studying English since primary school} : réponse exacte à la question 2 ; on évite le verbe d'état \eng{know} à la forme continue.
\item \eng{she will be preparing} et \eng{she will be working} : future continuous (\eng{will be} + V-ing) pour des actions en cours à un moment futur, avec les marqueurs \eng{this time next year} et \eng{in ten years}.
\item \eng{she will be living near her family} : \eng{live} accepte la forme continue ; attention, avec un verbe d'état comme \eng{know}, on dirait simplement \eng{she will know}.
\item Conclusion personnelle avec \eng{I am sure...} : prétérit modalisé \eng{will achieve} pour exprimer une prédiction.
\end{itemize}
\end{exoresolu}

\courssub{Bilan de l'activité}

\begin{aretenir}[Ce que cette activité t'a appris]
\begin{itemize}
\item Le \cle{present perfect continuous} (\eng{have/has been + V-ing}) exprime une action commencée dans le passé et qui dure encore ; il répond à la question \eng{How long...?} avec \eng{for} + durée ou \eng{since} + point de départ.
\item Les \cle{verbes d'état} (\eng{know}, \eng{like}, \eng{want}, \eng{believe}) refusent la forme continue : on utilise alors le present perfect simple (\eng{I have known her for years}).
\item Le \cle{future continuous} (\eng{will be + V-ing}) décrit une action en cours à un moment précis du futur, souvent introduite par \eng{this time next year} ou \eng{in ten years}.
\item En français, on utilise le présent avec « depuis » ; en anglais, c'est le present perfect continuous qui s'impose : \eng{I have been living here for ten years.}
\item À l'oral, soigne la prononciation de \eng{been} (« biine ») et la contraction naturelle \eng{I'll be}, \eng{she's been}.
\end{itemize}
\end{aretenir}

\begin{exercice}[Pour aller plus loin]
Corrige les erreurs dans les phrases suivantes, entendues pendant les lectures des portraits :
\begin{enumerate}
\item \eng{My classmate has living in Bamenda since ten years.}
\item \eng{She has been knowing her best friend for five years.}
\item \eng{In ten years, he will studying at the university.}
\item \eng{I am living in Douala since 2015.}
\end{enumerate}
\end{exercice}

\begin{corrige}[Exercice « Pour aller plus loin »]
\begin{enumerate}
\item \eng{My classmate has been living in Bamenda for ten years.} (ajout de \eng{been} ; \eng{for} + durée, \eng{since} + point de départ)
\item \eng{She has known her best friend for five years.} (\eng{know} est un verbe d'état : present perfect simple, pas de forme continue)
\item \eng{In ten years, he will be studying at the university.} (il manque \eng{be} au future continuous)
\item \eng{I have been living in Douala since 2015.} (avec \eng{since}, le présent français se traduit par le present perfect continuous)
\end{enumerate}
\end{corrige}

\newpage

\coursec{Méthodes et exercices résolus}

\courssub{Méthode 1 : Identifier le temps d'un verbe anglais}

Avant d'aborder les deux nouveaux temps de cette leçon, rappelons le panorama des temps étudiés depuis la classe de Seconde. Un verbe anglais se lit comme une formule : chaque auxiliaire apporte une information.

\begin{center}
\begin{tabularx}{\linewidth}{|l|X|X|}
\hline
\rowcolor{popblueL} \textbf{Temps} & \textbf{Forme} & \textbf{Exemple} \\
\hline
Present simple & V (V-s à la 3\textsuperscript{e} pers. sing.) & \eng{She works in Douala.} \\
\hline
Present continuous & am/is/are + V-ing & \eng{She is working now.} \\
\hline
Simple past & V-ed / 2\textsuperscript{e} colonne & \eng{She worked there in 2019.} \\
\hline
Past continuous & was/were + V-ing & \eng{She was working when I arrived.} \\
\hline
Present perfect simple & have/has + participe passé & \eng{She has finished her report.} \\
\hline
Past perfect & had + participe passé & \eng{She had left before noon.} \\
\hline
Future simple & will + V & \eng{She will work tomorrow.} \\
\hline
\rowcolor{popgreenL} Present perfect continuous & have/has + been + V-ing & \eng{She has been working since 8 a.m.} \\
\hline
\rowcolor{popgreenL} Future continuous & will + be + V-ing & \eng{She will be working at 10 a.m.} \\
\hline
\end{tabularx}
\end{center}

\begin{methode}[Identifier un temps verbal en 4 étapes]
Pour reconnaître le temps d'un verbe dans une phrase anglaise, procède méthodiquement :
\begin{enumerate}
\item \textbf{Repère le groupe verbal} (auxiliaire(s) + verbe principal). Exemple : \eng{has been working}.
\item \textbf{Analyse les auxiliaires} : \eng{have/has} = perfect ; \eng{be + -ing} = continuous (progressive) ; \eng{will} = future ; \eng{had} = past perfect.
\item \textbf{Combine les marques} : \eng{has been working} = have (perfect) + been + -ing (continuous) = \cle{present perfect continuous}.
\item \textbf{Vérifie avec les indices de temps} (time markers) : \eng{for, since, lately} $\rightarrow$ perfect continuous ; \eng{this time tomorrow, at 8 p.m. tonight} $\rightarrow$ future continuous ; \eng{yesterday, last year} $\rightarrow$ simple past.
\end{enumerate}
Réflexe : traduis mentalement la phrase ; si le français dit « cela fait ... que », pense au \cle{present perfect (continuous)}.
\end{methode}

\begin{exoresolu}[Identification des temps]
\textbf{Énoncé.} Identify the tense of each underlined verb and justify your answer.\\
1. \eng{My uncle \textbf{has been living} in Buea for ten years.}\\
2. \eng{This time next week, we \textbf{will be writing} our mock examination.}\\
3. \eng{She \textbf{visited} her grandmother in Bamenda last Christmas.}\\
4. \eng{They \textbf{have just finished} their homework.}
\tcblower
\textbf{Solution détaillée.}
\begin{enumerate}
\item \eng{has been living} : auxiliaire \eng{has} (perfect) + \eng{been} + \eng{-ing} (continuous) $\rightarrow$ \textbf{present perfect continuous}. L'indice \eng{for ten years} (durée qui continue jusqu'à maintenant) confirme. Traduction : « Mon oncle vit à Buea depuis dix ans. »
\item \eng{will be writing} : \eng{will} (future) + \eng{be} + \eng{-ing} (continuous) $\rightarrow$ \textbf{future continuous}. L'indice \eng{this time next week} désigne un moment précis du futur où l'action sera en cours : « À la même heure la semaine prochaine, nous serons en train de composer. »
\item \eng{visited} : verbe régulier au passé, sans auxiliaire $\rightarrow$ \textbf{simple past}. L'indice \eng{last Christmas} situe une action terminée et datée : « Elle a rendu visite à sa grand-mère à Bamenda à Noël dernier. »
\item \eng{have just finished} : \eng{have} + participe passé $\rightarrow$ \textbf{present perfect simple}. L'adverbe \eng{just} signale un résultat récent lié au présent : « Ils viennent de finir leurs devoirs. »
\end{enumerate}
\textbf{Conclusion.} Toujours croiser deux sources d'information : la \emph{forme} du groupe verbal et les \emph{indices temporels} de la phrase.
\end{exoresolu}

\courssub{Méthode 2 : Construire le present perfect continuous}

\begin{propriete}[Formation du present perfect continuous]
\begin{center}
\begin{tabular}{|l|l|l|}
\hline
\rowcolor{popblueL} \textbf{Forme} & \textbf{Structure} & \textbf{Exemple} \\
\hline
Affirmative & sujet + have/has + been + V-ing & \eng{They have been playing.} \\
\hline
Négative & sujet + have/has + not + been + V-ing & \eng{He has not been sleeping well.} \\
\hline
Interrogative & Have/Has + sujet + been + V-ing ? & \eng{Have you been waiting long?} \\
\hline
\end{tabular}
\end{center}
\eng{has} à la 3\textsuperscript{e} personne du singulier (he, she, it) ; \eng{have} partout ailleurs. Formes contractées : \eng{I've been, she's been, they haven't been}.
\end{propriete}

\begin{methode}[Former et employer le present perfect continuous]
\begin{enumerate}
\item \textbf{Formule} : \cle{subject + have/has + been + verb-ing}.
\item \textbf{Négation} : \eng{have/has + not + been + V-ing}. \eng{She has not been sleeping well.}
\item \textbf{Question} : inversion de l'auxiliaire : \eng{How long have you been waiting?}
\item \textbf{Emplois} : (a) action commencée dans le passé et \emph{qui continue} : \eng{I have been learning English since Form One.} ; (b) action récente dont le \emph{résultat est visible} : \eng{Your hands are dirty — have you been working in the garden?}
\item \textbf{Attention} : les verbes d'état (\eng{know, believe, love, understand, want}) ne prennent pas la forme continue : on dit \eng{I have known him for years}, jamais \emph{I have been knowing him}.
\end{enumerate}
\end{methode}

\begin{attention}[Orthographe du participe présent en -ing]
\begin{itemize}
\item Verbe en \eng{-e} muet : on supprime le \eng{e} : \eng{write $\rightarrow$ writing}, \eng{live $\rightarrow$ living} (mais \eng{be $\rightarrow$ being}, \eng{see $\rightarrow$ seeing}).
\item Consonne finale doublée après voyelle courte accentuée : \eng{sit $\rightarrow$ sitting}, \eng{run $\rightarrow$ running}, \eng{swim $\rightarrow$ swimming}.
\item \eng{-ie} devient \eng{-y} : \eng{lie $\rightarrow$ lying}, \eng{die $\rightarrow$ dying}.
\end{itemize}
\end{attention}

\begin{exoresolu}[Mise en forme au present perfect continuous]
\textbf{Énoncé.} Put the verbs in brackets into the present perfect continuous.\\
1. \eng{The students (revise) for the Probatoire all morning.}\\
2. \eng{It (rain) in Douala since six o'clock.}\\
3. \eng{How long (you / wait) for the bendskin?}\\
4. \eng{She (not / feel) well lately.}
\tcblower
\textbf{Solution détaillée.}
\begin{enumerate}
\item \eng{The students \textbf{have been revising} for the Probatoire all morning.} — sujet pluriel $\rightarrow$ \eng{have} ; \eng{all morning} indique une durée en cours. « Les élèves révisent pour le Probatoire depuis ce matin. »
\item \eng{It \textbf{has been raining} in Douala since six o'clock.} — sujet \eng{it} $\rightarrow$ \eng{has} ; \eng{since} + point de départ. « Il pleut à Douala depuis six heures. »
\item \eng{How long \textbf{have you been waiting} for the bendskin?} — interrogative : inversion \eng{have + sujet + been + V-ing}. « Depuis combien de temps attends-tu la moto-taxi ? »
\item \eng{She \textbf{has not been feeling} well lately.} — négation : \eng{has not been} + \eng{feeling} ; \eng{lately} (« dernièrement ») appelle ce temps. « Elle ne se sent pas bien ces derniers temps. »
\end{enumerate}
\textbf{Conclusion.} La charpente \eng{have/has been + V-ing} ne change jamais ; seul le choix \eng{have}/\eng{has} dépend du sujet.
\end{exoresolu}

\courssub{Méthode 3 : Choisir entre present perfect simple et present perfect continuous}

\begin{methode}[Simple ou continuous ? Le critère décisif]
\begin{enumerate}
\item Pose-toi la question : est-ce le \textbf{résultat / la quantité accomplie} qui compte, ou la \textbf{durée / l'activité en cours} ?
\item \textbf{Résultat, quantité, action achevée} $\rightarrow$ \cle{present perfect simple} : \eng{I have written three letters.} (« J'ai écrit trois lettres » — elles sont finies.)
\item \textbf{Durée, activité non terminée ou répétée} $\rightarrow$ \cle{present perfect continuous} : \eng{I have been writing letters all afternoon.} (« J'écris des lettres depuis tout l'après-midi » — l'activité compte, pas le nombre.)
\item Avec \eng{for} et \eng{since}, les deux sont souvent possibles, mais le continuous insiste sur la durée : \eng{He has worked here for five years} / \eng{He has been working here for five years} (plus vivant, plus actuel).
\item Verbes d'état $\rightarrow$ simple obligatoire : \eng{We have known each other since childhood.}
\end{enumerate}
\end{methode}

\begin{center}
\begin{tabularx}{\linewidth}{|X|X|}
\hline
\rowcolor{popblueL} \textbf{Present perfect simple (résultat)} & \textbf{Present perfect continuous (durée)} \\
\hline
\eng{I have read the novel.} (fini) & \eng{I have been reading the novel.} (pas fini) \\
\hline
\eng{She has painted the kitchen.} (c'est fait) & \eng{She has been painting the kitchen.} (elle est encore couverte de peinture) \\
\hline
\eng{How many pages have you written?} & \eng{How long have you been writing?} \\
\hline
\end{tabularx}
\end{center}

\begin{exoresolu}[Choix entre perfect simple et continuous]
\textbf{Énoncé.} Choose the correct form and justify your choice.\\
1. \eng{I (have read / have been reading) this novel; you can have it back now.}\\
2. \eng{She (has cooked / has been cooking) since four o'clock and the meal is not ready yet.}\\
3. \eng{They (have known / have been knowing) each other for twenty years.}\\
4. \eng{How many goals (has Eto'o scored / has Eto'o been scoring) this season?}
\tcblower
\textbf{Solution détaillée.}
\begin{enumerate}
\item \eng{I \textbf{have read} this novel; you can have it back now.} — le livre est terminé : c'est le \emph{résultat} qui compte $\rightarrow$ perfect simple. « J'ai lu ce roman ; tu peux le reprendre. »
\item \eng{She \textbf{has been cooking} since four o'clock and the meal is not ready yet.} — l'action \emph{continue} encore (\eng{not ready yet}) $\rightarrow$ continuous. « Elle cuisine depuis quatre heures et le repas n'est pas encore prêt. »
\item \eng{They \textbf{have known} each other for twenty years.} — \eng{know} est un verbe d'état : la forme continue est impossible. « Ils se connaissent depuis vingt ans. »
\item \eng{How many goals \textbf{has Eto'o scored} this season?} — \eng{how many} interroge sur une \emph{quantité accomplie} $\rightarrow$ perfect simple. « Combien de buts Eto'o a-t-il marqués cette saison ? »
\end{enumerate}
\textbf{Conclusion.} Nombre ou résultat = simple ; durée ou activité en cours = continuous ; verbe d'état = simple.
\end{exoresolu}

\courssub{Méthode 4 : Construire et employer le future continuous}

\begin{propriete}[Formation du future continuous]
\begin{center}
\begin{tabular}{|l|l|l|}
\hline
\rowcolor{popblueL} \textbf{Forme} & \textbf{Structure} & \textbf{Exemple} \\
\hline
Affirmative & sujet + will + be + V-ing & \eng{We will be travelling.} \\
\hline
Négative & sujet + will not (won't) + be + V-ing & \eng{She won't be coming.} \\
\hline
Interrogative & Will + sujet + be + V-ing ? & \eng{Will they be staying long?} \\
\hline
\end{tabular}
\end{center}
Forme invariable : \eng{will} ne change jamais, quelle que soit la personne.
\end{propriete}

\begin{methode}[Le future continuous : forme et emplois]
\begin{enumerate}
\item \textbf{Formule} : \cle{subject + will be + verb-ing}. Négation : \eng{will not (won't) be + V-ing}. Question : \eng{Will you be using the car tonight?}
\item \textbf{Emploi 1 — action en cours à un moment précis du futur} : \eng{At 10 a.m. tomorrow, we will be sitting our English paper.} (« Demain à 10 h, nous serons en train de composer. »)
\item \textbf{Emploi 2 — déroulement normal prévu} (sans intention particulière) : \eng{I will be travelling to Bamenda next week anyway, so I can take your parcel.}
\item \textbf{Emploi 3 — question polie sur les projets de quelqu'un} : \eng{Will you be coming to the family meeting on Sunday?}
\item \textbf{Indices typiques} : \eng{this time tomorrow, at this time next week, all day tomorrow}.
\item Ne pas confondre avec \eng{will + V} (décision ou prédiction simple) ni avec \eng{be going to} (intention ou projet).
\end{enumerate}
\end{methode}

\begin{aretenir}[Les trois futurs à ne pas confondre]
\begin{itemize}
\item \eng{I \textbf{will help} you.} — décision prise à l'instant / promesse. (« Je vais t'aider. »)
\item \eng{I \textbf{am going to help} her tomorrow.} — intention, projet déjà formé. (« J'ai l'intention de l'aider demain. »)
\item \eng{I \textbf{will be helping} the candidates all morning.} — action qui sera en cours à un moment du futur. (« Je serai en train d'aider les candidats toute la matinée. »)
\end{itemize}
\end{aretenir}

\begin{exoresolu}[Transformation vers le future continuous]
\textbf{Énoncé.} Rewrite each sentence using the future continuous, keeping the same meaning.\\
1. \eng{The family will travel to the village during the Christmas holidays.}\\
2. \eng{The teacher will mark our scripts tomorrow morning at eight.}\\
3. \eng{Will you use your phone this evening? I need to call my mother.}
\tcblower
\textbf{Solution détaillée.}
\begin{enumerate}
\item \eng{The family \textbf{will be travelling} to the village during the Christmas holidays.} — \eng{will be} + \eng{travelling} : l'action sera en cours pendant une période du futur. « La famille sera en train de voyager vers le village pendant les vacances de Noël. »
\item \eng{The teacher \textbf{will be marking} our scripts at eight tomorrow morning.} — moment précis du futur (\eng{at eight}) : le future continuous peint l'action en train de se dérouler. « Demain à huit heures, le professeur sera en train de corriger nos copies. »
\item \eng{\textbf{Will you be using} your phone this evening? I need to call my mother.} — question polie sur un projet : inversion \eng{Will + sujet + be + V-ing}. « Est-ce que tu utiliseras ton téléphone ce soir ? Je dois appeler ma mère. »
\end{enumerate}
\textbf{Conclusion.} La structure \eng{will be + V-ing} est mécanique ; c'est le \emph{sens} (action en cours dans le futur, question polie) qui guide son choix.
\end{exoresolu}

\courssub{Méthode 5 : Rédiger un court texte en mobilisant les temps étudiés}

\begin{methode}[Produire un paragraphe cohérent au Bac]
\begin{enumerate}
\item \textbf{Planifie} : situation présente (present perfect continuous pour les activités en cours), bilan (present perfect simple pour les résultats), projection (future continuous pour les actions futures en cours).
\item \textbf{Place les connecteurs temporels} : \eng{for, since, lately, so far, this time next month}.
\item \textbf{Varie les temps sans les mélanger au hasard} : chaque temps doit être justifié par le sens.
\item \textbf{Relis} en vérifiant : accord \eng{have/has}, participe passé des verbes irréguliers, \eng{-ing} après \eng{been} et après \eng{will be}.
\end{enumerate}
Structure modèle : \eng{I have been ... for/since ... / I have already ... / This time next year, I will be ...}
\end{methode}

\begin{exoresolu}[Composition guidée]
\textbf{Énoncé.} Write a short paragraph (about 60 words) about your life as a candidate for the Baccalauréat, using at least once the present perfect continuous, the present perfect simple and the future continuous.
\tcblower
\textbf{Solution détaillée.}\\
Étape 1 — activité en cours : \eng{I have been preparing for the Baccalauréat since September.} (present perfect continuous + \eng{since}).\\
Étape 2 — résultats déjà obtenus : \eng{I have already revised all my grammar notes and I have written three practice essays.} (present perfect simple, quantités accomplies).\\
Étape 3 — projection : \eng{This time next month, I will be sitting the written papers at Lycée de Nkol-Eton.} (future continuous, moment précis du futur).\\
\textbf{Réponse rédigée :}\\
\eng{I have been preparing for the Baccalauréat since September, and I have worked harder than ever this year. I have already revised all my grammar notes and written three practice essays. My teachers have been encouraging us a lot lately. This time next month, I will be sitting the written papers, and I am confident that I will succeed.}\\
« Je prépare le Baccalauréat depuis septembre, et j'ai travaillé plus que jamais cette année. J'ai déjà révisé toutes mes notes de grammaire et rédigé trois dissertations d'entraînement. Nos professeurs nous encouragent beaucoup ces derniers temps. À la même époque le mois prochain, je serai en train de composer les épreuves écrites, et je suis sûr de réussir. »\\
\textbf{Conclusion.} Trois temps, trois fonctions : durée présente, bilan accompli, action future en cours — la logique temporelle fait la cohérence du texte.
\end{exoresolu}

\begin{attention}[Les 5 erreurs qui coûtent des points au Bac]
\begin{enumerate}
\item \textbf{Calque du « depuis » français} : \emph{I live in Yaoundé since 2020} est faux. Le français utilise le présent, l'anglais exige le perfect : \eng{I \textbf{have been living} in Yaoundé since 2020.} De même : \eng{How long \textbf{have you been waiting}?} et non \emph{How long do you wait?}
\item \textbf{Oubli de \eng{been}} : \emph{She has working hard} est impossible. La forme complète est \eng{She \textbf{has been working} hard.} Trois éléments obligatoires : \eng{has} + \eng{been} + \eng{working}.
\item \textbf{Verbe d'état à la forme continue} : \emph{I have been knowing him for years} est une faute grave. Avec \eng{know, believe, love, want, understand}, utilise le perfect simple : \eng{I \textbf{have known} him for years.}
\item \textbf{Future sans \eng{be}} : \emph{This time tomorrow, I will writing my exam} est incorrect. Il faut \eng{I \textbf{will be writing} my exam.} Le \eng{be} est indispensable avant le \eng{-ing}.
\item \textbf{Confusion \eng{for}/\eng{since}} : \eng{for} + durée (\eng{for ten years}), \eng{since} + point de départ (\eng{since 2015, since Christmas}). Dire \emph{since ten years} est un calque du français « depuis dix ans ».
\end{enumerate}
\end{attention}

\begin{exercice}[Pour t'entraîner]
\begin{enumerate}
\item Put the verbs in brackets into the correct tense (present perfect simple, present perfect continuous or future continuous) :\\
a) \eng{We (wait) for the results for two weeks.}\\
b) \eng{My sister (just / pass) her driving test.}\\
c) \eng{This time on Friday, the candidates (write) the English paper.}\\
d) \eng{How long (you / know) your best friend?}
\item Correct the mistakes in these sentences :\\
a) \eng{I am living in Bamenda since 2018.}\\
b) \eng{She has been know him for years.}\\
c) \eng{Tomorrow at nine, we will writing the test.}
\item Write five sentences about your family using the three tenses of this lesson.
\end{enumerate}
\end{exercice}

\begin{corrige}[Exercice « Pour t'entraîner »]
\begin{enumerate}
\item a) \eng{We \textbf{have been waiting} for the results for two weeks.} (durée en cours) ; b) \eng{My sister \textbf{has just passed} her driving test.} (résultat récent, \eng{just}) ; c) \eng{This time on Friday, the candidates \textbf{will be writing} the English paper.} (moment précis du futur) ; d) \eng{How long \textbf{have you known} your best friend?} (\eng{know} = verbe d'état, perfect simple).
\item a) \eng{I \textbf{have been living} in Bamenda since 2018.} ; b) \eng{She \textbf{has known} him for years.} ; c) \eng{Tomorrow at nine, we \textbf{will be writing} the test.}
\item Exemples : \eng{My father has been working at the same company for fifteen years. / We have just visited our grandparents in the village. / This time next Sunday, the whole family will be celebrating my sister's birthday.}
\end{enumerate}
\end{corrige}

\newpage

\coursec{Exercices}

\courssub{Je vérifie mes connaissances}

\begin{exercice}[QCM : la bonne forme verbale]
Pour chaque phrase, choisis la bonne réponse (a, b ou c).
\begin{enumerate}
\item How long \ldots\ English?\\
a) do you learn \quad b) have you been learning \quad c) are you learning
\item This time tomorrow, the candidates \ldots\ in the examination hall.\\
a) will sit \quad b) will be sitting \quad c) sit
\item She is exhausted because she \ldots\ the house all day.\\
a) has cleaned \quad b) has been cleaning \quad c) cleans
\item I \ldots\ three essays this week. (résultat : ils sont finis)\\
a) have written \quad b) have been writing \quad c) am writing
\item Next July, my sister \ldots\ at the University of Buea.\\
a) will studying \quad b) will be studying \quad c) is studying
\end{enumerate}
\end{exercice}

\begin{exercice}[Vrai ou faux ? Justifie]
Dis si chaque affirmation est vraie ou fausse, puis justifie ta réponse par la règle de grammaire.
\begin{enumerate}
\item On peut dire \eng{I have been knowing her since 2019.}
\item Le present perfect continuous s'emploie pour insister sur la durée d'une action commencée dans le passé et qui continue encore.
\item \eng{We are living here since ten years} est une phrase correcte en anglais.
\item Le future continuous se forme avec \eng{will be} + verbe en \eng{-ing}.
\item Avec \eng{since} et \eng{for}, les verbes d'état comme \eng{know}, \eng{believe} et \eng{love} se mettent normalement au present perfect simple, pas au continuous.
\end{enumerate}
\end{exercice}

\begin{exercice}[Vocabulaire : les marqueurs de temps]
Relie chaque expression de temps à l'emploi qui lui correspond, puis utilise chacune dans une phrase de ton invention au temps qui convient.
\begin{center}
\begin{tabularx}{\linewidth}{|X|X|}
\hline
\rowcolor{popblueL} Expression de temps & Emploi \\
\hline
\eng{since 2015} & point de départ dans le passé \\
\hline
\eng{for ten years} & durée \\
\hline
\eng{this time tomorrow} & action en cours à un moment précis du futur \\
\hline
\eng{all morning} & durée continue jusqu'à maintenant \\
\hline
\eng{by next June} & action achevée avant une limite future \\
\hline
\end{tabularx}
\end{center}
\end{exercice}

\begin{exercice}[Conjugaison à compléter]
Complète les phrases avec le verbe entre parenthèses au present perfect simple ou au present perfect continuous, selon le sens.
\begin{enumerate}
\item I \ldots\ (write) three essays this week. They are ready.
\item I \ldots\ (write) essays all morning and I am tired.
\item My father \ldots\ (work) at the same company for twenty years.
\item Look at the mud on your shoes! \ldots\ you \ldots\ (play) football?
\item We \ldots\ (not / finish) our project yet.
\item She \ldots\ (learn) to drive since January, but she still has no licence.
\end{enumerate}
\end{exercice}

\courssub{Je m'entraîne}

\begin{exercice}[Traduction français vers anglais]
Traduis les phrases suivantes en anglais, en choisissant le temps correct.
\begin{enumerate}
\item Mon frère travaille à Douala depuis cinq ans. À la même heure demain, il sera en train de conduire vers Yaoundé.
\item Nous révisons pour le Baccalauréat depuis le mois de septembre.
\item Elle est fatiguée parce qu'elle vend des légumes au marché depuis six heures du matin.
\item Cette fois-ci la semaine prochaine, nous serons en train de passer l'épreuve d'anglais.
\item Depuis combien de temps attends-tu le bendskin ?
\end{enumerate}
\end{exercice}

\begin{exercice}[Transformation de phrases]
Réécris chaque phrase en utilisant le temps indiqué entre parenthèses, sans changer le sens.
\begin{enumerate}
\item He started working here in 2015 and still works here. (present perfect continuous)
\item The rain began at six o'clock and it is still raining. (present perfect continuous)
\item Tomorrow at 8 a.m., the students will sit for the English paper. (future continuous)
\item She began teaching in Bamenda three years ago and she is still there. (present perfect continuous)
\item Don't phone me at seven tonight; I will have dinner with my family. (future continuous)
\end{enumerate}
\end{exercice}

\begin{exercice}[Texte à trous]
Complète le texte suivant avec les verbes entre parenthèses au temps qui convient : present perfect simple, present perfect continuous ou future continuous.
\begin{textebox}[A letter from Yaoundé]
Dear Paul,

Greetings from Yaoundé! I \ldots\ (live) here for six months now, and I \ldots\ already \ldots\ (make) many friends. Life is busy: I \ldots\ (prepare) for the Baccalauréat since September, and I \ldots\ (attend) extra classes every Saturday. My parents \ldots\ (support) me a lot — they \ldots\ just \ldots\ (buy) me a new dictionary.

This time next month, we \ldots\ (write) our final examinations. I am nervous, but I know that while I \ldots\ (sit) in the examination hall, my whole family \ldots\ (think) of me.

Write soon and tell me what you \ldots\ (do) these days.

Your friend,
Alain
\end{textebox}
\end{exercice}

\begin{exercice}[Corrige les erreurs]
Chaque phrase contient une erreur typique d'élève francophone. Souligne l'erreur et réécris la phrase correctement.
\begin{enumerate}
\item I have been know her since 2019.
\item She will studying in Buea next year.
\item We are living here since ten years.
\item He has been working here since five years ago.
\item How long time have you been waiting for me?
\item They have been married since twenty years.
\end{enumerate}
\end{exercice}

\courssub{Je me prépare au Bac}

\begin{exercice}[Compréhension et langue — type Bac]
Lis le texte suivant, puis réponds aux questions.

\begin{textebox}[Waiting for the results]
For the past three weeks, thousands of candidates in Cameroon have been waiting anxiously for the Baccalauréat results. In Yaoundé, nineteen-year-old Nadine has been checking the notice board of her school every morning. “I have been dreaming about this moment since Form Five,” she says. “I have worked very hard, and now I just want to know.”

Her younger brother Eric, who will sit for the examination next year, has already started preparing. “This time next year, I will be writing my own papers,” he explains. Their mother, a trader at Mokolo market, has been saving money for Nadine's university fees for five years. “Education is the best gift,” she says with a smile.

While families have been waiting, teachers have been marking scripts day and night. Soon, the results will be published, and while some students will be celebrating, others will be planning to try again.
\end{textebox}

\textbf{A. Comprehension questions} — Answer in complete sentences.
\begin{enumerate}
\item How long have the candidates been waiting for their results?
\item What has Nadine been doing every morning?
\item What will Eric be doing this time next year?
\item Why has Nadine's mother been saving money?
\item Say whether the following statement is true or false and justify with a quotation from the text: “The teachers finished marking the scripts long ago.”
\end{enumerate}

\textbf{B. Language work}
\begin{enumerate}
\item Find in the text two verbs in the present perfect continuous and two verbs in the future continuous. Copy them and underline the auxiliary.
\item Rewrite the following sentence in the future continuous: “The results will be published soon.”
\item Complete with the correct tense: “Nadine \ldots\ (wait) for her results for three weeks now.”
\end{enumerate}
\end{exercice}

\begin{exercice}[Traduction — type Bac]
Traduis le passage suivant en anglais. Soigne particulièrement le choix des temps verbaux.

« Depuis deux ans, ma sœur prépare son concours d'entrée à l'ENAM. Elle étudie tous les soirs et elle a déjà lu plus de vingt livres. Cette fois-ci le mois prochain, elle sera en train de composer dans la salle d'examen. Toute la famille pense à elle et nous sommes sûrs qu'elle réussira. »
\end{exercice}

\begin{exercice}[Composition — type Bac]
Write a composition of about \textbf{80 to 100 words} on the following topic:

\emph{Describe what your family has been doing recently and what you will be doing after the Baccalauréat.}

Consignes :
\begin{itemize}
\item Use the present perfect continuous at least twice and the future continuous at least twice.
\item Organise your text in two short paragraphs (recent activities / future plans).
\item Use time expressions such as \eng{since}, \eng{for}, \eng{this time next month}, \eng{all week}.
\end{itemize}
\end{exercice}

\newpage

\coursec{Corrigés des exercices}

\begin{corrige}[Exercice 1 — QCM : la bonne forme verbale]
\begin{enumerate}
\item \textbf{b) have you been learning} — \eng{How long have you been learning English?} (« Depuis combien de temps apprends-tu l'anglais ? ») La question porte sur la durée d'une action commencée dans le passé et qui continue : present perfect continuous, imposé par \eng{how long}.
\item \textbf{b) will be sitting} — \eng{This time tomorrow, the candidates will be sitting in the examination hall.} (« À la même heure demain, les candidats seront assis dans la salle d'examen. ») Le marqueur \eng{this time tomorrow} exige le future continuous : action en cours à un moment précis du futur.
\item \textbf{b) has been cleaning} — \eng{She is exhausted because she has been cleaning the house all day.} (« Elle est épuisée parce qu'elle a nettoyé la maison toute la journée. ») On explique un résultat présent (la fatigue) par une activité longue et récente (\eng{all day}) : present perfect continuous.
\item \textbf{a) have written} — \eng{I have written three essays this week.} (« J'ai écrit trois dissertations cette semaine. ») On insiste sur le résultat achevé et quantifié (trois essais terminés) : present perfect simple, pas continuous.
\item \textbf{b) will be studying} — \eng{Next July, my sister will be studying at the University of Buea.} (« En juillet prochain, ma sœur étudiera à l'université de Buea. ») Future continuous : \eng{will be} + verbe en \eng{-ing}. La réponse a) est fautive : \eng{will studying} n'existe pas (il manque \eng{be}).
\end{enumerate}
\end{corrige}

\begin{corrige}[Exercice 2 — Vrai ou faux ? Justifie]
\begin{enumerate}
\item \textbf{Faux.} \eng{Know} est un verbe d'état : il ne se met pas à la forme continue. On dit \eng{I have known her since 2019.} (« Je la connais depuis 2019. »), jamais \eng{I have been knowing her}.
\item \textbf{Vrai.} C'est l'emploi principal du present perfect continuous : insister sur la durée d'une action commencée dans le passé et qui se poursuit au moment où l'on parle (\eng{I have been waiting for an hour.} — « J'attends depuis une heure. »).
\item \textbf{Faux.} Deux erreurs : avec une durée qui continue jusqu'à maintenant, il faut le present perfect (continuous) et non le présent simple ; de plus, \eng{since} s'emploie avec un point de départ, pas avec une durée. Correction : \eng{We have been living here for ten years.} (« Nous habitons ici depuis dix ans. »)
\item \textbf{Vrai.} Formation du future continuous : sujet + \eng{will be} + verbe en \eng{-ing} (\eng{They will be travelling.} — « Ils seront en train de voyager. »).
\item \textbf{Vrai.} Les verbes d'état (\eng{know, believe, love, understand, want}...) n'admettent normalement pas la forme continue : \eng{I have loved her for years.} et non \eng{I have been loving her}.
\end{enumerate}
\end{corrige}

\begin{corrige}[Exercice 3 — Vocabulaire : les marqueurs de temps]
Correspondances (déjà données dans le tableau) et phrases modèles :
\begin{itemize}
\item \eng{since 2015} (point de départ) : \eng{My uncle has been living in Bamenda since 2015.} (« Mon oncle vit à Bamenda depuis 2015. »)
\item \eng{for ten years} (durée) : \eng{They have known each other for ten years.} (« Ils se connaissent depuis dix ans. » — verbe d'état, donc present perfect simple.)
\item \eng{this time tomorrow} (action en cours dans le futur) : \eng{This time tomorrow, we will be sitting for the English paper.} (« À la même heure demain, nous serons en train de composer l'épreuve d'anglais. »)
\item \eng{all morning} (durée continue jusqu'à maintenant) : \eng{It has been raining all morning.} (« Il pleut depuis ce matin. »)
\item \eng{by next June} (limite future) : \eng{By next June, I will have finished my studies.} (« D'ici juin prochain, j'aurai terminé mes études. » — future perfect : action achevée avant une date future.)
\end{itemize}
\end{corrige}

\begin{corrige}[Exercice 4 — Conjugaison à compléter]
\begin{enumerate}
\item \eng{I \textbf{have written} three essays this week. They are ready.} — résultat achevé et quantifié : present perfect simple.
\item \eng{I \textbf{have been writing} essays all morning and I am tired.} — durée (\eng{all morning}) et résultat visible (la fatigue) : present perfect continuous.
\item \eng{My father \textbf{has been working} at the same company for twenty years.} — action qui dure (\eng{for twenty years}) ; le simple \eng{has worked} est aussi correct, mais le continuous insiste mieux sur la continuité.
\item \eng{\textbf{Have} you \textbf{been playing} football?} — traces de boue visibles : activité récente dont le résultat se voit encore.
\item \eng{We \textbf{have not finished} our project yet.} — \eng{yet} à la forme négative : present perfect simple (résultat attendu, non obtenu).
\item \eng{She \textbf{has been learning} to drive since January, but she still has no licence.} — apprentissage en cours depuis un point de départ (\eng{since January}) : continuous.
\end{enumerate}
\end{corrige}

\begin{corrige}[Exercice 5 — Traduction français vers anglais]
\begin{enumerate}
\item \eng{My brother has been working in Douala for five years. This time tomorrow, he will be driving to Yaoundé.} (« depuis cinq ans » = \eng{for five years} ; « sera en train de conduire » = future continuous.)
\item \eng{We have been revising for the Baccalauréat since September.} (point de départ : \eng{since}.)
\item \eng{She is tired because she has been selling vegetables at the market since six o'clock in the morning.}
\item \eng{This time next week, we will be sitting for the English paper.} (\eng{sit for an exam/a paper} = passer une épreuve.)
\item \eng{How long have you been waiting for the bendskin?} (Attention : \eng{wait for} + complément ; ne pas oublier la préposition \eng{for}.)
\end{enumerate}
\end{corrige}

\begin{corrige}[Exercice 6 — Transformation de phrases]
\begin{enumerate}
\item \eng{He has been working here since 2015.}
\item \eng{It has been raining since six o'clock.}
\item \eng{Tomorrow at 8 a.m., the students will be sitting for the English paper.}
\item \eng{She has been teaching in Bamenda for three years.}
\item \eng{Don't phone me at seven tonight; I will be having dinner with my family.}
\end{enumerate}
Points de vigilance : \eng{since} + point de départ (2015, six o'clock), \eng{for} + durée (three years) ; au future continuous, on garde \eng{will be} + \eng{-ing} (\eng{will be sitting}, \eng{will be having} — \eng{have} perd son \eng{e} final devant \eng{-ing}).
\end{corrige}

\begin{corrige}[Exercice 7 — Texte à trous]
Texte complété :

\eng{Dear Paul,}

\eng{Greetings from Yaoundé! I \textbf{have been living} here for six months now, and I \textbf{have} already \textbf{made} many friends. Life is busy: I \textbf{have been preparing} for the Baccalauréat since September, and I \textbf{attend} extra classes every Saturday. My parents \textbf{have supported} me a lot — they \textbf{have} just \textbf{bought} me a new dictionary.}

\eng{This time next month, we \textbf{will be writing} our final examinations. I am nervous, but I know that while I \textbf{am sitting} in the examination hall, my whole family \textbf{will be thinking} of me.}

\eng{Write soon and tell me what you \textbf{have been doing} these days.}

\eng{Your friend,}
\eng{Alain}

Justifications :
\begin{itemize}
\item \eng{have been living} : durée (\eng{for six months}) qui continue ; \eng{have already made} : résultat acquis ; \eng{have been preparing} : activité en cours depuis septembre ; \eng{attend} : habitude présente (\eng{every Saturday}) ; \eng{have supported} et \eng{have just bought} : actions récentes avec résultat présent (\eng{just}).
\item \eng{will be writing} : \eng{this time next month} impose le future continuous ; \eng{am sitting} : après \eng{while}, le présent continuous remplace le futur (jamais \eng{will} après \eng{while} ou \eng{when} au sens temporel) ; \eng{will be thinking} : action simultanée dans le futur ; \eng{have been doing} : activités récentes jusqu'à maintenant (\eng{these days}).
\end{itemize}
\end{corrige}

\begin{corrige}[Exercice 8 — Corrige les erreurs]
\begin{enumerate}
\item Phrase fautive : \eng{I have been knowing her since 2019.} Correction : \eng{I \textbf{have known} her since 2019.} (verbe d'état : pas de forme continue.)
\item Phrase fautive : \eng{She will studying in Buea next year.} Correction : \eng{She \textbf{will be studying} in Buea next year.} (future continuous : \eng{will be} + \eng{-ing}.)
\item Phrase fautive : \eng{We are living here since ten years.} Correction : \eng{We \textbf{have been living} here \textbf{for} ten years.}
\item Phrase fautive : \eng{He has been working here since five years ago.} Correction : \eng{He has been working here \textbf{for five years}.} (\eng{since} + point de départ : \eng{since 2019} ; \eng{for} + durée ; on ne combine jamais \eng{since} et \eng{ago}.)
\item Phrase fautive : \eng{How long time have you been waiting for me?} Correction : \eng{\textbf{How long} have you been waiting for me?} (\eng{how long} suffit ; ne pas traduire « combien de temps » mot à mot.)
\item Phrase fautive : \eng{They have been married since twenty years.} Correction : \eng{They have been married \textbf{for} twenty years.}
\end{enumerate}
\end{corrige}

\begin{corrige}[Exercice 9 — Compréhension et langue — type Bac]
\textbf{Barème indicatif (sur 20)} : A. Comprehension : 5 $\times$ 2 pts = 10 pts ; B. Language work : 4 + 3 + 3 = 10 pts.

\textbf{A. Comprehension questions — réponses modèles :}
\begin{enumerate}
\item \eng{The candidates have been waiting for their results for the past three weeks.}
\item \eng{Every morning, Nadine has been checking the notice board of her school.}
\item \eng{This time next year, Eric will be writing his own examination papers.}
\item \eng{Nadine's mother has been saving money to pay for Nadine's university fees.}
\item \textbf{False.} \eng{The statement is false because the text says teachers “have been marking scripts day and night”, which means the marking is still going on.}
\end{enumerate}

\textbf{B. Language work :}
\begin{enumerate}
\item Present perfect continuous : \eng{have \textbf{been waiting}}, \eng{has \textbf{been checking}} (également : \eng{has been dreaming}, \eng{has been saving}, \eng{have been marking}). Future continuous : \eng{will \textbf{be writing}}, \eng{will \textbf{be celebrating}} (également : \eng{will be planning}). L'auxiliaire à souligner est \eng{have/has} (suivi de \eng{been}) pour le premier temps, et \eng{will} pour le second.
\item Transformation attendue : \eng{Soon, they will be publishing the results.} (passage à la voix active au future continuous, forme naturelle). À éviter : \eng{The results will be being published soon.} — cette forme passive du future continuous, bien que théoriquement construite, est lourde et quasiment jamais employée en anglais réel ; elle serait pénalisée au Bac.
\item \eng{Nadine \textbf{has been waiting} for her results for three weeks now.}
\end{enumerate}

\textbf{Points de vigilance :} répondre par des phrases complètes ; ne pas recopier le texte sans l'adapter (changer les pronoms et les possessifs) ; dans la justification du vrai/faux, citer le texte entre guillemets anglais.
\end{corrige}

\begin{corrige}[Exercice 10 — Traduction — type Bac]
\textbf{Barème indicatif (sur 10)} : choix des temps 4 pts ; vocabulaire 3 pts ; orthographe et syntaxe 3 pts.

\textbf{Traduction modèle :}

\eng{For two years, my sister has been preparing for the ENAM entrance examination. She studies every evening and she has already read more than twenty books. This time next month, she will be sitting for the examination in the examination hall. The whole family is thinking of her and we are sure that she will succeed.}

\textbf{Points de vigilance :}
\begin{itemize}
\item « Depuis deux ans » + action qui continue : \eng{for two years} + present perfect continuous (\eng{has been preparing}), jamais le présent simple.
\item « elle a déjà lu » : present perfect simple avec \eng{already} (\eng{has already read}) — participe passé irrégulier \eng{read} (prononcé « rèd »).
\item « sera en train de composer » : future continuous \eng{will be sitting for} ; « composer une épreuve » se dit \eng{sit for an examination}, pas \eng{compose}.
\item « toute la famille pense à elle » : présent simple ou continuous (\eng{is thinking of her}) ; attention à \eng{think of} et non \eng{think to}.
\end{itemize}
\end{corrige}

\begin{corrige}[Exercice 11 — Composition — type Bac]
\textbf{Barème indicatif (sur 20)} : respect des consignes et des temps imposés 6 pts ; organisation en deux paragraphes 4 pts ; vocabulaire et marqueurs de temps 4 pts ; correction grammaticale et orthographe 6 pts.

\textbf{Composition modèle (environ 95 mots) :}

\eng{Recently, my family has been very busy. My mother has been preparing my younger brother for the Common Entrance Examination, and my father has been working extra hours to pay our school fees. As for me, I have been revising for the Baccalauréat all week, and I feel quite confident.}

\eng{After the examination, life will change. This time next month, I will be resting at my grandmother's village near Bafoussam. In September, I will be studying at university, and my brother will be starting secondary school. We are all looking forward to this new life.}

\textbf{Points de vigilance :}
\begin{itemize}
\item Vérifier la présence d'au moins deux present perfect continuous (\eng{has been preparing}, \eng{has been working}, \eng{have been revising}) et deux future continuous (\eng{will be resting}, \eng{will be studying}, \eng{will be starting}).
\item Employer les marqueurs demandés : \eng{all week}, \eng{this time next month}.
\item Éviter les calques du français : « je serai en train d'étudier » se dit simplement \eng{I will be studying} ; jamais \eng{I will be to study}.
\item Compter les mots et rester entre 80 et 100 ; soigner la présentation en deux paragraphes nets.
\end{itemize}
\end{corrige}

\newpage

\coursec{Fiche bilan}

\begin{popfigure}
\begin{tikzpicture}[mindmap, grow cyclic, every node/.style={concept}, root concept/.append style={concept color=popblue, font=\bfseries, text=white}, level 1/.append style={level distance=3.4cm, sibling angle=51}, level 2/.append style={level distance=2.2cm, sibling angle=40}]
\node[root concept, align=center]{Tenses Review\\ Tle A}
  child[concept color=poporange]{ node[align=center]{Simple present\\ habits, facts}
    child{ node[concept color=poporangeL, text=popdark]{She works in Douala} }}
  child[concept color=popgreen]{ node[align=center]{Present continuous\\ now, around now}
    child{ node[concept color=popgreenL, text=popdark]{She is revising now} }}
  child[concept color=poppurple]{ node[align=center]{Past simple vs\\ past continuous}
    child{ node[concept color=poppurpleL, text=popdark]{While I was walking, it rained} }}
  child[concept color=popgold]{ node[align=center]{Present perfect\\ simple}
    child{ node[concept color=popgoldL, text=popdark]{I have finished} }}
  child[concept color=poppink]{ node{Present perfect continuous}
    child{ node[concept color=poppinkL, text=popdark]{have been + V-ing} }}
  child[concept color=popblue!70]{ node[align=center]{Future simple\\ will / going to}
    child{ node[concept color=popblueL, text=popdark]{I will help you} }}
  child[concept color=popgreen!80!popblue]{ node{Future continuous}
    child{ node[concept color=popgreenL, text=popdark]{will be + V-ing} }};
\end{tikzpicture}
\legende{Carte mentale : les temps anglais révisés en Tle A et les deux nouveaux temps continus.}
\end{popfigure}

\begin{bilan}
\textbf{1. Lexique clé (marqueurs de temps)}

\begin{center}
\begin{tabularx}{\linewidth}{|X|X|}\hline
\rowcolor{popblueL} English & Français \\ \hline
for three hours / for years & depuis trois heures / depuis des années (durée) \\ \hline
since 2019 / since morning & depuis 2019 / depuis le matin (point de départ) \\ \hline
all day, all week, lately, recently & toute la journée, toute la semaine, dernièrement \\ \hline
this time tomorrow, at 8 p.m. tonight & demain à la même heure, ce soir à 20 h \\ \hline
already, just, yet, ever, never & déjà, juste, encore (nég.), jamais (interr.), jamais \\ \hline
\end{tabularx}
\end{center}

\textbf{2. Règles à connaître par cœur}

\begin{center}
\begin{tabularx}{\linewidth}{|l|X|X|}\hline
\rowcolor{popblueL} Temps & Formation & Emploi \\ \hline
Present perfect continuous & \cle{have/has been + V-ing} & action commencée dans le passé, encore en cours ou aux effets visibles ; insiste sur la durée \\ \hline
Future continuous & \cle{will be + V-ing} & action en cours à un moment précis du futur ; action prévue, naturelle \\ \hline
Present perfect simple & have/has + participe passé & résultat, bilan, quantité accomplie \\ \hline
Past simple / past continuous & V-ed / was-were + V-ing & action passée terminée / action longue interrompue \\ \hline
\end{tabularx}
\end{center}

\textbf{3. Contrastes décisifs}
\begin{itemize}
\item \eng{I have read three chapters.} (résultat : trois chapitres lus) vs \eng{I have been reading all morning.} (durée, activité en cours).
\item \eng{She has lived in Buea since 2020.} = \eng{She has been living in Buea since 2020.} (les deux possibles avec \eng{live}, \eng{work}, \eng{study}).
\item Verbes d'état (\eng{know, like, believe, want}) : jamais de forme continue : \eng{I have known her for years.} (jamais \emph{have been knowing}).
\item Futur : \eng{I will see her at 5.} (décision/promesse) ; \eng{I will be seeing her at 5.} (c'est prévu, cela arrivera dans le cours normal des choses).
\end{itemize}

\textbf{4. Méthodes express}
\begin{itemize}
\item Repérer le marqueur (\eng{for}, \eng{since}, \eng{this time tomorrow}) avant de choisir le temps.
\item Question clé : durée en cours $\rightarrow$ perfect continuous ; résultat $\rightarrow$ perfect simple ; moment futur précis $\rightarrow$ future continuous.
\item Transformation : \eng{He started working at 8 and he is still working.} $\rightarrow$ \eng{He has been working since 8 o'clock.}
\end{itemize}
\end{bilan}

\begin{aretenir}[Checklist avant le Bac]
\begin{itemize}
\item[$\square$] Je sais former le present perfect continuous : sujet + have/has + been + V-ing.
\item[$\square$] Je sais former le future continuous : sujet + will be + V-ing.
\item[$\square$] Je distingue \eng{for} (durée) et \eng{since} (point de départ).
\item[$\square$] Je choisis entre perfect simple (résultat) et perfect continuous (durée).
\item[$\square$] Je n'utilise jamais la forme continue avec les verbes d'état (\eng{know, love, believe}).
\item[$\square$] Je n'emploie jamais « I am knowing » ni « I am having been ».
\item[$\square$] Je n'oublie pas le \eng{-ing} après \eng{been} et après \eng{will be}.
\item[$\square$] Je n'ajoute pas de \eng{-s} à la 3\up{e} personne après \eng{will} : \eng{she will be working}.
\item[$\square$] Je relie les marqueurs de temps au bon temps verbal.
\item[$\square$] Je sais transformer une phrase du past simple au present perfect continuous.
\item[$\square$] Je sais employer le future continuous pour parler d'un programme prévu : \eng{This time tomorrow, we will be sitting the exam.}
\item[$\square$] Je relis ma production pour vérifier l'accord have/has avec le sujet.
\end{itemize}
\end{aretenir}

\findecours

\end{document}
